\documentclass[11pt]{article}

    \usepackage[breakable]{tcolorbox}
    \usepackage{parskip} % Stop auto-indenting (to mimic markdown behaviour)
    

    % Basic figure setup, for now with no caption control since it's done
    % automatically by Pandoc (which extracts ![](path) syntax from Markdown).
    \usepackage{graphicx}
    % Keep aspect ratio if custom image width or height is specified
    \setkeys{Gin}{keepaspectratio}
    % Maintain compatibility with old templates. Remove in nbconvert 6.0
    \let\Oldincludegraphics\includegraphics
    % Ensure that by default, figures have no caption (until we provide a
    % proper Figure object with a Caption API and a way to capture that
    % in the conversion process - todo).
    \usepackage{caption}
    \DeclareCaptionFormat{nocaption}{}
    \captionsetup{format=nocaption,aboveskip=0pt,belowskip=0pt}

    \usepackage{float}
    \floatplacement{figure}{H} % forces figures to be placed at the correct location
    \usepackage{xcolor} % Allow colors to be defined
    \usepackage{enumerate} % Needed for markdown enumerations to work
    \usepackage{geometry} % Used to adjust the document margins
    \usepackage{amsmath} % Equations
    \usepackage{amssymb} % Equations
    \usepackage{textcomp} % defines textquotesingle
    % Hack from http://tex.stackexchange.com/a/47451/13684:
    \AtBeginDocument{%
        \def\PYZsq{\textquotesingle}% Upright quotes in Pygmentized code
    }
    \usepackage{upquote} % Upright quotes for verbatim code
    \usepackage{eurosym} % defines \euro

    \usepackage{iftex}
    \ifPDFTeX
        \usepackage[T1]{fontenc}
        \IfFileExists{alphabeta.sty}{
              \usepackage{alphabeta}
          }{
              \usepackage[mathletters]{ucs}
              \usepackage[utf8x]{inputenc}
          }
    \else
        \usepackage{fontspec}
        \usepackage{unicode-math}
    \fi

    \usepackage{fancyvrb} % verbatim replacement that allows latex
    \usepackage{grffile} % extends the file name processing of package graphics
                         % to support a larger range
    \makeatletter % fix for old versions of grffile with XeLaTeX
    \@ifpackagelater{grffile}{2019/11/01}
    {
      % Do nothing on new versions
    }
    {
      \def\Gread@@xetex#1{%
        \IfFileExists{"\Gin@base".bb}%
        {\Gread@eps{\Gin@base.bb}}%
        {\Gread@@xetex@aux#1}%
      }
    }
    \makeatother
    \usepackage[Export]{adjustbox} % Used to constrain images to a maximum size
    \adjustboxset{max size={0.9\linewidth}{0.9\paperheight}}

    % The hyperref package gives us a pdf with properly built
    % internal navigation ('pdf bookmarks' for the table of contents,
    % internal cross-reference links, web links for URLs, etc.)
    \usepackage{hyperref}
    % The default LaTeX title has an obnoxious amount of whitespace. By default,
    % titling removes some of it. It also provides customization options.
    \usepackage{titling}
    \usepackage{longtable} % longtable support required by pandoc >1.10
    \usepackage{booktabs}  % table support for pandoc > 1.12.2
    \usepackage{array}     % table support for pandoc >= 2.11.3
    \usepackage{calc}      % table minipage width calculation for pandoc >= 2.11.1
    \usepackage[inline]{enumitem} % IRkernel/repr support (it uses the enumerate* environment)
    \usepackage[normalem]{ulem} % ulem is needed to support strikethroughs (\sout)
                                % normalem makes italics be italics, not underlines
    \usepackage{soul}      % strikethrough (\st) support for pandoc >= 3.0.0
    \usepackage{mathrsfs}
    

    
    % Colors for the hyperref package
    \definecolor{urlcolor}{rgb}{0,.145,.698}
    \definecolor{linkcolor}{rgb}{.71,0.21,0.01}
    \definecolor{citecolor}{rgb}{.12,.54,.11}

    % ANSI colors
    \definecolor{ansi-black}{HTML}{3E424D}
    \definecolor{ansi-black-intense}{HTML}{282C36}
    \definecolor{ansi-red}{HTML}{E75C58}
    \definecolor{ansi-red-intense}{HTML}{B22B31}
    \definecolor{ansi-green}{HTML}{00A250}
    \definecolor{ansi-green-intense}{HTML}{007427}
    \definecolor{ansi-yellow}{HTML}{DDB62B}
    \definecolor{ansi-yellow-intense}{HTML}{B27D12}
    \definecolor{ansi-blue}{HTML}{208FFB}
    \definecolor{ansi-blue-intense}{HTML}{0065CA}
    \definecolor{ansi-magenta}{HTML}{D160C4}
    \definecolor{ansi-magenta-intense}{HTML}{A03196}
    \definecolor{ansi-cyan}{HTML}{60C6C8}
    \definecolor{ansi-cyan-intense}{HTML}{258F8F}
    \definecolor{ansi-white}{HTML}{C5C1B4}
    \definecolor{ansi-white-intense}{HTML}{A1A6B2}
    \definecolor{ansi-default-inverse-fg}{HTML}{FFFFFF}
    \definecolor{ansi-default-inverse-bg}{HTML}{000000}

    % common color for the border for error outputs.
    \definecolor{outerrorbackground}{HTML}{FFDFDF}

    % commands and environments needed by pandoc snippets
    % extracted from the output of `pandoc -s`
    \providecommand{\tightlist}{%
      \setlength{\itemsep}{0pt}\setlength{\parskip}{0pt}}
    \DefineVerbatimEnvironment{Highlighting}{Verbatim}{commandchars=\\\{\}}
    % Add ',fontsize=\small' for more characters per line
    \newenvironment{Shaded}{}{}
    \newcommand{\KeywordTok}[1]{\textcolor[rgb]{0.00,0.44,0.13}{\textbf{{#1}}}}
    \newcommand{\DataTypeTok}[1]{\textcolor[rgb]{0.56,0.13,0.00}{{#1}}}
    \newcommand{\DecValTok}[1]{\textcolor[rgb]{0.25,0.63,0.44}{{#1}}}
    \newcommand{\BaseNTok}[1]{\textcolor[rgb]{0.25,0.63,0.44}{{#1}}}
    \newcommand{\FloatTok}[1]{\textcolor[rgb]{0.25,0.63,0.44}{{#1}}}
    \newcommand{\CharTok}[1]{\textcolor[rgb]{0.25,0.44,0.63}{{#1}}}
    \newcommand{\StringTok}[1]{\textcolor[rgb]{0.25,0.44,0.63}{{#1}}}
    \newcommand{\CommentTok}[1]{\textcolor[rgb]{0.38,0.63,0.69}{\textit{{#1}}}}
    \newcommand{\OtherTok}[1]{\textcolor[rgb]{0.00,0.44,0.13}{{#1}}}
    \newcommand{\AlertTok}[1]{\textcolor[rgb]{1.00,0.00,0.00}{\textbf{{#1}}}}
    \newcommand{\FunctionTok}[1]{\textcolor[rgb]{0.02,0.16,0.49}{{#1}}}
    \newcommand{\RegionMarkerTok}[1]{{#1}}
    \newcommand{\ErrorTok}[1]{\textcolor[rgb]{1.00,0.00,0.00}{\textbf{{#1}}}}
    \newcommand{\NormalTok}[1]{{#1}}

    % Additional commands for more recent versions of Pandoc
    \newcommand{\ConstantTok}[1]{\textcolor[rgb]{0.53,0.00,0.00}{{#1}}}
    \newcommand{\SpecialCharTok}[1]{\textcolor[rgb]{0.25,0.44,0.63}{{#1}}}
    \newcommand{\VerbatimStringTok}[1]{\textcolor[rgb]{0.25,0.44,0.63}{{#1}}}
    \newcommand{\SpecialStringTok}[1]{\textcolor[rgb]{0.73,0.40,0.53}{{#1}}}
    \newcommand{\ImportTok}[1]{{#1}}
    \newcommand{\DocumentationTok}[1]{\textcolor[rgb]{0.73,0.13,0.13}{\textit{{#1}}}}
    \newcommand{\AnnotationTok}[1]{\textcolor[rgb]{0.38,0.63,0.69}{\textbf{\textit{{#1}}}}}
    \newcommand{\CommentVarTok}[1]{\textcolor[rgb]{0.38,0.63,0.69}{\textbf{\textit{{#1}}}}}
    \newcommand{\VariableTok}[1]{\textcolor[rgb]{0.10,0.09,0.49}{{#1}}}
    \newcommand{\ControlFlowTok}[1]{\textcolor[rgb]{0.00,0.44,0.13}{\textbf{{#1}}}}
    \newcommand{\OperatorTok}[1]{\textcolor[rgb]{0.40,0.40,0.40}{{#1}}}
    \newcommand{\BuiltInTok}[1]{{#1}}
    \newcommand{\ExtensionTok}[1]{{#1}}
    \newcommand{\PreprocessorTok}[1]{\textcolor[rgb]{0.74,0.48,0.00}{{#1}}}
    \newcommand{\AttributeTok}[1]{\textcolor[rgb]{0.49,0.56,0.16}{{#1}}}
    \newcommand{\InformationTok}[1]{\textcolor[rgb]{0.38,0.63,0.69}{\textbf{\textit{{#1}}}}}
    \newcommand{\WarningTok}[1]{\textcolor[rgb]{0.38,0.63,0.69}{\textbf{\textit{{#1}}}}}
    \makeatletter
    \newsavebox\pandoc@box
    \newcommand*\pandocbounded[1]{%
      \sbox\pandoc@box{#1}%
      % scaling factors for width and height
      \Gscale@div\@tempa\textheight{\dimexpr\ht\pandoc@box+\dp\pandoc@box\relax}%
      \Gscale@div\@tempb\linewidth{\wd\pandoc@box}%
      % select the smaller of both
      \ifdim\@tempb\p@<\@tempa\p@
        \let\@tempa\@tempb
      \fi
      % scaling accordingly (\@tempa < 1)
      \ifdim\@tempa\p@<\p@
        \scalebox{\@tempa}{\usebox\pandoc@box}%
      % scaling not needed, use as it is
      \else
        \usebox{\pandoc@box}%
      \fi
    }
    \makeatother

    % Define a nice break command that doesn't care if a line doesn't already
    % exist.
    \def\br{\hspace*{\fill} \\* }
    % Math Jax compatibility definitions
    \def\gt{>}
    \def\lt{<}
    \let\Oldtex\TeX
    \let\Oldlatex\LaTeX
    \renewcommand{\TeX}{\textrm{\Oldtex}}
    \renewcommand{\LaTeX}{\textrm{\Oldlatex}}
    % Document parameters
    % Document title
    \title{RGR1\_A-2}
    
    
    
    
    
    
    
% Pygments definitions
\makeatletter
\def\PY@reset{\let\PY@it=\relax \let\PY@bf=\relax%
    \let\PY@ul=\relax \let\PY@tc=\relax%
    \let\PY@bc=\relax \let\PY@ff=\relax}
\def\PY@tok#1{\csname PY@tok@#1\endcsname}
\def\PY@toks#1+{\ifx\relax#1\empty\else%
    \PY@tok{#1}\expandafter\PY@toks\fi}
\def\PY@do#1{\PY@bc{\PY@tc{\PY@ul{%
    \PY@it{\PY@bf{\PY@ff{#1}}}}}}}
\def\PY#1#2{\PY@reset\PY@toks#1+\relax+\PY@do{#2}}

\@namedef{PY@tok@w}{\def\PY@tc##1{\textcolor[rgb]{0.73,0.73,0.73}{##1}}}
\@namedef{PY@tok@c}{\let\PY@it=\textit\def\PY@tc##1{\textcolor[rgb]{0.24,0.48,0.48}{##1}}}
\@namedef{PY@tok@cp}{\def\PY@tc##1{\textcolor[rgb]{0.61,0.40,0.00}{##1}}}
\@namedef{PY@tok@k}{\let\PY@bf=\textbf\def\PY@tc##1{\textcolor[rgb]{0.00,0.50,0.00}{##1}}}
\@namedef{PY@tok@kp}{\def\PY@tc##1{\textcolor[rgb]{0.00,0.50,0.00}{##1}}}
\@namedef{PY@tok@kt}{\def\PY@tc##1{\textcolor[rgb]{0.69,0.00,0.25}{##1}}}
\@namedef{PY@tok@o}{\def\PY@tc##1{\textcolor[rgb]{0.40,0.40,0.40}{##1}}}
\@namedef{PY@tok@ow}{\let\PY@bf=\textbf\def\PY@tc##1{\textcolor[rgb]{0.67,0.13,1.00}{##1}}}
\@namedef{PY@tok@nb}{\def\PY@tc##1{\textcolor[rgb]{0.00,0.50,0.00}{##1}}}
\@namedef{PY@tok@nf}{\def\PY@tc##1{\textcolor[rgb]{0.00,0.00,1.00}{##1}}}
\@namedef{PY@tok@nc}{\let\PY@bf=\textbf\def\PY@tc##1{\textcolor[rgb]{0.00,0.00,1.00}{##1}}}
\@namedef{PY@tok@nn}{\let\PY@bf=\textbf\def\PY@tc##1{\textcolor[rgb]{0.00,0.00,1.00}{##1}}}
\@namedef{PY@tok@ne}{\let\PY@bf=\textbf\def\PY@tc##1{\textcolor[rgb]{0.80,0.25,0.22}{##1}}}
\@namedef{PY@tok@nv}{\def\PY@tc##1{\textcolor[rgb]{0.10,0.09,0.49}{##1}}}
\@namedef{PY@tok@no}{\def\PY@tc##1{\textcolor[rgb]{0.53,0.00,0.00}{##1}}}
\@namedef{PY@tok@nl}{\def\PY@tc##1{\textcolor[rgb]{0.46,0.46,0.00}{##1}}}
\@namedef{PY@tok@ni}{\let\PY@bf=\textbf\def\PY@tc##1{\textcolor[rgb]{0.44,0.44,0.44}{##1}}}
\@namedef{PY@tok@na}{\def\PY@tc##1{\textcolor[rgb]{0.41,0.47,0.13}{##1}}}
\@namedef{PY@tok@nt}{\let\PY@bf=\textbf\def\PY@tc##1{\textcolor[rgb]{0.00,0.50,0.00}{##1}}}
\@namedef{PY@tok@nd}{\def\PY@tc##1{\textcolor[rgb]{0.67,0.13,1.00}{##1}}}
\@namedef{PY@tok@s}{\def\PY@tc##1{\textcolor[rgb]{0.73,0.13,0.13}{##1}}}
\@namedef{PY@tok@sd}{\let\PY@it=\textit\def\PY@tc##1{\textcolor[rgb]{0.73,0.13,0.13}{##1}}}
\@namedef{PY@tok@si}{\let\PY@bf=\textbf\def\PY@tc##1{\textcolor[rgb]{0.64,0.35,0.47}{##1}}}
\@namedef{PY@tok@se}{\let\PY@bf=\textbf\def\PY@tc##1{\textcolor[rgb]{0.67,0.36,0.12}{##1}}}
\@namedef{PY@tok@sr}{\def\PY@tc##1{\textcolor[rgb]{0.64,0.35,0.47}{##1}}}
\@namedef{PY@tok@ss}{\def\PY@tc##1{\textcolor[rgb]{0.10,0.09,0.49}{##1}}}
\@namedef{PY@tok@sx}{\def\PY@tc##1{\textcolor[rgb]{0.00,0.50,0.00}{##1}}}
\@namedef{PY@tok@m}{\def\PY@tc##1{\textcolor[rgb]{0.40,0.40,0.40}{##1}}}
\@namedef{PY@tok@gh}{\let\PY@bf=\textbf\def\PY@tc##1{\textcolor[rgb]{0.00,0.00,0.50}{##1}}}
\@namedef{PY@tok@gu}{\let\PY@bf=\textbf\def\PY@tc##1{\textcolor[rgb]{0.50,0.00,0.50}{##1}}}
\@namedef{PY@tok@gd}{\def\PY@tc##1{\textcolor[rgb]{0.63,0.00,0.00}{##1}}}
\@namedef{PY@tok@gi}{\def\PY@tc##1{\textcolor[rgb]{0.00,0.52,0.00}{##1}}}
\@namedef{PY@tok@gr}{\def\PY@tc##1{\textcolor[rgb]{0.89,0.00,0.00}{##1}}}
\@namedef{PY@tok@ge}{\let\PY@it=\textit}
\@namedef{PY@tok@gs}{\let\PY@bf=\textbf}
\@namedef{PY@tok@ges}{\let\PY@bf=\textbf\let\PY@it=\textit}
\@namedef{PY@tok@gp}{\let\PY@bf=\textbf\def\PY@tc##1{\textcolor[rgb]{0.00,0.00,0.50}{##1}}}
\@namedef{PY@tok@go}{\def\PY@tc##1{\textcolor[rgb]{0.44,0.44,0.44}{##1}}}
\@namedef{PY@tok@gt}{\def\PY@tc##1{\textcolor[rgb]{0.00,0.27,0.87}{##1}}}
\@namedef{PY@tok@err}{\def\PY@bc##1{{\setlength{\fboxsep}{\string -\fboxrule}\fcolorbox[rgb]{1.00,0.00,0.00}{1,1,1}{\strut ##1}}}}
\@namedef{PY@tok@kc}{\let\PY@bf=\textbf\def\PY@tc##1{\textcolor[rgb]{0.00,0.50,0.00}{##1}}}
\@namedef{PY@tok@kd}{\let\PY@bf=\textbf\def\PY@tc##1{\textcolor[rgb]{0.00,0.50,0.00}{##1}}}
\@namedef{PY@tok@kn}{\let\PY@bf=\textbf\def\PY@tc##1{\textcolor[rgb]{0.00,0.50,0.00}{##1}}}
\@namedef{PY@tok@kr}{\let\PY@bf=\textbf\def\PY@tc##1{\textcolor[rgb]{0.00,0.50,0.00}{##1}}}
\@namedef{PY@tok@bp}{\def\PY@tc##1{\textcolor[rgb]{0.00,0.50,0.00}{##1}}}
\@namedef{PY@tok@fm}{\def\PY@tc##1{\textcolor[rgb]{0.00,0.00,1.00}{##1}}}
\@namedef{PY@tok@vc}{\def\PY@tc##1{\textcolor[rgb]{0.10,0.09,0.49}{##1}}}
\@namedef{PY@tok@vg}{\def\PY@tc##1{\textcolor[rgb]{0.10,0.09,0.49}{##1}}}
\@namedef{PY@tok@vi}{\def\PY@tc##1{\textcolor[rgb]{0.10,0.09,0.49}{##1}}}
\@namedef{PY@tok@vm}{\def\PY@tc##1{\textcolor[rgb]{0.10,0.09,0.49}{##1}}}
\@namedef{PY@tok@sa}{\def\PY@tc##1{\textcolor[rgb]{0.73,0.13,0.13}{##1}}}
\@namedef{PY@tok@sb}{\def\PY@tc##1{\textcolor[rgb]{0.73,0.13,0.13}{##1}}}
\@namedef{PY@tok@sc}{\def\PY@tc##1{\textcolor[rgb]{0.73,0.13,0.13}{##1}}}
\@namedef{PY@tok@dl}{\def\PY@tc##1{\textcolor[rgb]{0.73,0.13,0.13}{##1}}}
\@namedef{PY@tok@s2}{\def\PY@tc##1{\textcolor[rgb]{0.73,0.13,0.13}{##1}}}
\@namedef{PY@tok@sh}{\def\PY@tc##1{\textcolor[rgb]{0.73,0.13,0.13}{##1}}}
\@namedef{PY@tok@s1}{\def\PY@tc##1{\textcolor[rgb]{0.73,0.13,0.13}{##1}}}
\@namedef{PY@tok@mb}{\def\PY@tc##1{\textcolor[rgb]{0.40,0.40,0.40}{##1}}}
\@namedef{PY@tok@mf}{\def\PY@tc##1{\textcolor[rgb]{0.40,0.40,0.40}{##1}}}
\@namedef{PY@tok@mh}{\def\PY@tc##1{\textcolor[rgb]{0.40,0.40,0.40}{##1}}}
\@namedef{PY@tok@mi}{\def\PY@tc##1{\textcolor[rgb]{0.40,0.40,0.40}{##1}}}
\@namedef{PY@tok@il}{\def\PY@tc##1{\textcolor[rgb]{0.40,0.40,0.40}{##1}}}
\@namedef{PY@tok@mo}{\def\PY@tc##1{\textcolor[rgb]{0.40,0.40,0.40}{##1}}}
\@namedef{PY@tok@ch}{\let\PY@it=\textit\def\PY@tc##1{\textcolor[rgb]{0.24,0.48,0.48}{##1}}}
\@namedef{PY@tok@cm}{\let\PY@it=\textit\def\PY@tc##1{\textcolor[rgb]{0.24,0.48,0.48}{##1}}}
\@namedef{PY@tok@cpf}{\let\PY@it=\textit\def\PY@tc##1{\textcolor[rgb]{0.24,0.48,0.48}{##1}}}
\@namedef{PY@tok@c1}{\let\PY@it=\textit\def\PY@tc##1{\textcolor[rgb]{0.24,0.48,0.48}{##1}}}
\@namedef{PY@tok@cs}{\let\PY@it=\textit\def\PY@tc##1{\textcolor[rgb]{0.24,0.48,0.48}{##1}}}

\def\PYZbs{\char`\\}
\def\PYZus{\char`\_}
\def\PYZob{\char`\{}
\def\PYZcb{\char`\}}
\def\PYZca{\char`\^}
\def\PYZam{\char`\&}
\def\PYZlt{\char`\<}
\def\PYZgt{\char`\>}
\def\PYZsh{\char`\#}
\def\PYZpc{\char`\%}
\def\PYZdl{\char`\$}
\def\PYZhy{\char`\-}
\def\PYZsq{\char`\'}
\def\PYZdq{\char`\"}
\def\PYZti{\char`\~}
% for compatibility with earlier versions
\def\PYZat{@}
\def\PYZlb{[}
\def\PYZrb{]}
\makeatother


    % For linebreaks inside Verbatim environment from package fancyvrb.
    \makeatletter
        \newbox\Wrappedcontinuationbox
        \newbox\Wrappedvisiblespacebox
        \newcommand*\Wrappedvisiblespace {\textcolor{red}{\textvisiblespace}}
        \newcommand*\Wrappedcontinuationsymbol {\textcolor{red}{\llap{\tiny$\m@th\hookrightarrow$}}}
        \newcommand*\Wrappedcontinuationindent {3ex }
        \newcommand*\Wrappedafterbreak {\kern\Wrappedcontinuationindent\copy\Wrappedcontinuationbox}
        % Take advantage of the already applied Pygments mark-up to insert
        % potential linebreaks for TeX processing.
        %        {, <, #, %, $, ' and ": go to next line.
        %        _, }, ^, &, >, - and ~: stay at end of broken line.
        % Use of \textquotesingle for straight quote.
        \newcommand*\Wrappedbreaksatspecials {%
            \def\PYGZus{\discretionary{\char`\_}{\Wrappedafterbreak}{\char`\_}}%
            \def\PYGZob{\discretionary{}{\Wrappedafterbreak\char`\{}{\char`\{}}%
            \def\PYGZcb{\discretionary{\char`\}}{\Wrappedafterbreak}{\char`\}}}%
            \def\PYGZca{\discretionary{\char`\^}{\Wrappedafterbreak}{\char`\^}}%
            \def\PYGZam{\discretionary{\char`\&}{\Wrappedafterbreak}{\char`\&}}%
            \def\PYGZlt{\discretionary{}{\Wrappedafterbreak\char`\<}{\char`\<}}%
            \def\PYGZgt{\discretionary{\char`\>}{\Wrappedafterbreak}{\char`\>}}%
            \def\PYGZsh{\discretionary{}{\Wrappedafterbreak\char`\#}{\char`\#}}%
            \def\PYGZpc{\discretionary{}{\Wrappedafterbreak\char`\%}{\char`\%}}%
            \def\PYGZdl{\discretionary{}{\Wrappedafterbreak\char`\$}{\char`\$}}%
            \def\PYGZhy{\discretionary{\char`\-}{\Wrappedafterbreak}{\char`\-}}%
            \def\PYGZsq{\discretionary{}{\Wrappedafterbreak\textquotesingle}{\textquotesingle}}%
            \def\PYGZdq{\discretionary{}{\Wrappedafterbreak\char`\"}{\char`\"}}%
            \def\PYGZti{\discretionary{\char`\~}{\Wrappedafterbreak}{\char`\~}}%
        }
        % Some characters . , ; ? ! / are not pygmentized.
        % This macro makes them "active" and they will insert potential linebreaks
        \newcommand*\Wrappedbreaksatpunct {%
            \lccode`\~`\.\lowercase{\def~}{\discretionary{\hbox{\char`\.}}{\Wrappedafterbreak}{\hbox{\char`\.}}}%
            \lccode`\~`\,\lowercase{\def~}{\discretionary{\hbox{\char`\,}}{\Wrappedafterbreak}{\hbox{\char`\,}}}%
            \lccode`\~`\;\lowercase{\def~}{\discretionary{\hbox{\char`\;}}{\Wrappedafterbreak}{\hbox{\char`\;}}}%
            \lccode`\~`\:\lowercase{\def~}{\discretionary{\hbox{\char`\:}}{\Wrappedafterbreak}{\hbox{\char`\:}}}%
            \lccode`\~`\?\lowercase{\def~}{\discretionary{\hbox{\char`\?}}{\Wrappedafterbreak}{\hbox{\char`\?}}}%
            \lccode`\~`\!\lowercase{\def~}{\discretionary{\hbox{\char`\!}}{\Wrappedafterbreak}{\hbox{\char`\!}}}%
            \lccode`\~`\/\lowercase{\def~}{\discretionary{\hbox{\char`\/}}{\Wrappedafterbreak}{\hbox{\char`\/}}}%
            \catcode`\.\active
            \catcode`\,\active
            \catcode`\;\active
            \catcode`\:\active
            \catcode`\?\active
            \catcode`\!\active
            \catcode`\/\active
            \lccode`\~`\~
        }
    \makeatother

    \let\OriginalVerbatim=\Verbatim
    \makeatletter
    \renewcommand{\Verbatim}[1][1]{%
        %\parskip\z@skip
        \sbox\Wrappedcontinuationbox {\Wrappedcontinuationsymbol}%
        \sbox\Wrappedvisiblespacebox {\FV@SetupFont\Wrappedvisiblespace}%
        \def\FancyVerbFormatLine ##1{\hsize\linewidth
            \vtop{\raggedright\hyphenpenalty\z@\exhyphenpenalty\z@
                \doublehyphendemerits\z@\finalhyphendemerits\z@
                \strut ##1\strut}%
        }%
        % If the linebreak is at a space, the latter will be displayed as visible
        % space at end of first line, and a continuation symbol starts next line.
        % Stretch/shrink are however usually zero for typewriter font.
        \def\FV@Space {%
            \nobreak\hskip\z@ plus\fontdimen3\font minus\fontdimen4\font
            \discretionary{\copy\Wrappedvisiblespacebox}{\Wrappedafterbreak}
            {\kern\fontdimen2\font}%
        }%

        % Allow breaks at special characters using \PYG... macros.
        \Wrappedbreaksatspecials
        % Breaks at punctuation characters . , ; ? ! and / need catcode=\active
        \OriginalVerbatim[#1,codes*=\Wrappedbreaksatpunct]%
    }
    \makeatother

    % Exact colors from NB
    \definecolor{incolor}{HTML}{303F9F}
    \definecolor{outcolor}{HTML}{D84315}
    \definecolor{cellborder}{HTML}{CFCFCF}
    \definecolor{cellbackground}{HTML}{F7F7F7}

    % prompt
    \makeatletter
    \newcommand{\boxspacing}{\kern\kvtcb@left@rule\kern\kvtcb@boxsep}
    \makeatother
    \newcommand{\prompt}[4]{
        {\ttfamily\llap{{\color{#2}[#3]:\hspace{3pt}#4}}\vspace{-\baselineskip}}
    }
    

    
    % Prevent overflowing lines due to hard-to-break entities
    \sloppy
    % Setup hyperref package
    \hypersetup{
      breaklinks=true,  % so long urls are correctly broken across lines
      colorlinks=true,
      urlcolor=urlcolor,
      linkcolor=linkcolor,
      citecolor=citecolor,
      }
    % Slightly bigger margins than the latex defaults
    
    \geometry{verbose,tmargin=1in,bmargin=1in,lmargin=1in,rmargin=1in}
    
    
\usepackage{fontspec}
\setmainfont{Arial} 
\setsansfont{Arial}
\setmonofont{Courier New}
\usepackage[russian]{babel}
\begin{document}
    
    \maketitle
    
    

    
    \subsection{Расчетно-графическая работа по математической статистике.
Вариант A-2. Выполнил Григорьев Даниил, поток 21.1,
P3216}\label{ux440ux430ux441ux447ux435ux442ux43dux43e-ux433ux440ux430ux444ux438ux447ux435ux441ux43aux430ux44f-ux440ux430ux431ux43eux442ux430-ux43fux43e-ux43cux430ux442ux435ux43cux430ux442ux438ux447ux435ux441ux43aux43eux439-ux441ux442ux430ux442ux438ux441ux442ux438ux43aux435.-ux432ux430ux440ux438ux430ux43dux442-a-2.-ux432ux44bux43fux43eux43bux43dux438ux43b-ux433ux440ux438ux433ux43eux440ux44cux435ux432-ux434ux430ux43dux438ux438ux43b-ux43fux43eux442ux43eux43a-21.1-p3216}

Группа A, единицы измерения --- мс, объем выборки \(n = 200\).

    \subsubsection{1. Первичное описание
выборки}\label{ux43fux435ux440ux432ux438ux447ux43dux43eux435-ux43eux43fux438ux441ux430ux43dux438ux435-ux432ux44bux431ux43eux440ux43aux438}

    \begin{tcolorbox}[breakable, size=fbox, boxrule=1pt, pad at break*=1mm,colback=cellbackground, colframe=cellborder]
\prompt{In}{incolor}{2}{\boxspacing}
\begin{Verbatim}[commandchars=\\\{\}]
\PY{o}{!}pip\PY{+w}{ }install\PY{+w}{ }scikit\PYZhy{}learn
\end{Verbatim}
\end{tcolorbox}

    \begin{Verbatim}[commandchars=\\\{\}]
Defaulting to user installation because normal site-packages is not writeable
Requirement already satisfied: scikit-learn in
C:\textbackslash{}Users\textbackslash{}danie\textbackslash{}AppData\textbackslash{}Roaming\textbackslash{}Python\textbackslash{}Python314\textbackslash{}site-packages (1.8.0)
Requirement already satisfied: numpy>=1.24.1 in
C:\textbackslash{}Users\textbackslash{}danie\textbackslash{}AppData\textbackslash{}Roaming\textbackslash{}Python\textbackslash{}Python314\textbackslash{}site-packages (from scikit-
learn) (2.4.2)
Requirement already satisfied: scipy>=1.10.0 in
C:\textbackslash{}Users\textbackslash{}danie\textbackslash{}AppData\textbackslash{}Roaming\textbackslash{}Python\textbackslash{}Python314\textbackslash{}site-packages (from scikit-
learn) (1.17.1)
Requirement already satisfied: joblib>=1.3.0 in
C:\textbackslash{}Users\textbackslash{}danie\textbackslash{}AppData\textbackslash{}Roaming\textbackslash{}Python\textbackslash{}Python314\textbackslash{}site-packages (from scikit-
learn) (1.5.3)
Requirement already satisfied: threadpoolctl>=3.2.0 in
C:\textbackslash{}Users\textbackslash{}danie\textbackslash{}AppData\textbackslash{}Roaming\textbackslash{}Python\textbackslash{}Python314\textbackslash{}site-packages (from scikit-
learn) (3.6.0)
    \end{Verbatim}

    \begin{tcolorbox}[breakable, size=fbox, boxrule=1pt, pad at break*=1mm,colback=cellbackground, colframe=cellborder]
\prompt{In}{incolor}{3}{\boxspacing}
\begin{Verbatim}[commandchars=\\\{\}]
\PY{k+kn}{import}\PY{+w}{ }\PY{n+nn}{pandas}\PY{+w}{ }\PY{k}{as}\PY{+w}{ }\PY{n+nn}{pd}
\PY{k+kn}{import}\PY{+w}{ }\PY{n+nn}{numpy}\PY{+w}{ }\PY{k}{as}\PY{+w}{ }\PY{n+nn}{np}
\PY{k+kn}{import}\PY{+w}{ }\PY{n+nn}{matplotlib}\PY{n+nn}{.}\PY{n+nn}{pyplot}\PY{+w}{ }\PY{k}{as}\PY{+w}{ }\PY{n+nn}{plt}
\PY{k+kn}{import}\PY{+w}{ }\PY{n+nn}{scipy}\PY{n+nn}{.}\PY{n+nn}{stats}\PY{+w}{ }\PY{k}{as}\PY{+w}{ }\PY{n+nn}{stats}
\PY{k+kn}{from}\PY{+w}{ }\PY{n+nn}{scipy}\PY{+w}{ }\PY{k+kn}{import} \PY{n}{optimize}
\PY{k+kn}{import}\PY{+w}{ }\PY{n+nn}{math}
\PY{k+kn}{from}\PY{+w}{ }\PY{n+nn}{pathlib}\PY{+w}{ }\PY{k+kn}{import} \PY{n}{Path}
\PY{k+kn}{from}\PY{+w}{ }\PY{n+nn}{IPython}\PY{n+nn}{.}\PY{n+nn}{display}\PY{+w}{ }\PY{k+kn}{import} \PY{n}{display}\PY{p}{,} \PY{n}{Latex}\PY{p}{,} \PY{n}{Math}
\PY{k+kn}{from}\PY{+w}{ }\PY{n+nn}{sklearn}\PY{n+nn}{.}\PY{n+nn}{cluster}\PY{+w}{ }\PY{k+kn}{import} \PY{n}{KMeans}
\PY{k+kn}{from}\PY{+w}{ }\PY{n+nn}{six}\PY{+w}{ }\PY{k+kn}{import} \PY{n}{string\PYZus{}types}
\PY{k+kn}{from}\PY{+w}{ }\PY{n+nn}{statsmodels}\PY{n+nn}{.}\PY{n+nn}{distributions}\PY{n+nn}{.}\PY{n+nn}{empirical\PYZus{}distribution}\PY{+w}{ }\PY{k+kn}{import} \PY{n}{ECDF}
\end{Verbatim}
\end{tcolorbox}

    \begin{tcolorbox}[breakable, size=fbox, boxrule=1pt, pad at break*=1mm,colback=cellbackground, colframe=cellborder]
\prompt{In}{incolor}{4}{\boxspacing}
\begin{Verbatim}[commandchars=\\\{\}]
\PY{c+c1}{\PYZsh{} 1. Загрузка данных}
\PY{n}{possible\PYZus{}paths} \PY{o}{=} \PY{p}{[}
    \PY{n}{Path}\PY{p}{(}\PY{l+s+s1}{\PYZsq{}}\PY{l+s+s1}{RGR1\PYZus{}A\PYZhy{}2\PYZus{}X1\PYZhy{}X4.csv}\PY{l+s+s1}{\PYZsq{}}\PY{p}{)}\PY{p}{,}
    \PY{n}{Path}\PY{p}{(}\PY{l+s+s1}{\PYZsq{}}\PY{l+s+s1}{/mnt/data/RGR1\PYZus{}A\PYZhy{}2\PYZus{}X1\PYZhy{}X4.csv}\PY{l+s+s1}{\PYZsq{}}\PY{p}{)}\PY{p}{,}
\PY{p}{]}

\PY{k}{for} \PY{n}{path} \PY{o+ow}{in} \PY{n}{possible\PYZus{}paths}\PY{p}{:}
    \PY{k}{if} \PY{n}{path}\PY{o}{.}\PY{n}{exists}\PY{p}{(}\PY{p}{)}\PY{p}{:}
        \PY{n}{df} \PY{o}{=} \PY{n}{pd}\PY{o}{.}\PY{n}{read\PYZus{}csv}\PY{p}{(}\PY{n}{path}\PY{p}{)}
        \PY{k}{break}
\PY{k}{else}\PY{p}{:}
    \PY{k}{raise} \PY{n+ne}{FileNotFoundError}\PY{p}{(}\PY{l+s+s1}{\PYZsq{}}\PY{l+s+s1}{Файл RGR1\PYZus{}A\PYZhy{}2\PYZus{}X1\PYZhy{}X4.csv не найден}\PY{l+s+s1}{\PYZsq{}}\PY{p}{)}

\PY{c+c1}{\PYZsh{} 2. Объем выборки}
\PY{n}{n} \PY{o}{=} \PY{n+nb}{len}\PY{p}{(}\PY{n}{df}\PY{p}{)}
\PY{n+nb}{print}\PY{p}{(}\PY{l+s+sa}{f}\PY{l+s+s2}{\PYZdq{}}\PY{l+s+s2}{Объем выборки n = }\PY{l+s+si}{\PYZob{}}\PY{n}{n}\PY{l+s+si}{\PYZcb{}}\PY{l+s+s2}{\PYZdq{}}\PY{p}{)}

\PY{c+c1}{\PYZsh{} 3. Разделение на переменные}
\PY{n}{X1} \PY{o}{=} \PY{n}{df}\PY{p}{[}\PY{l+s+s1}{\PYZsq{}}\PY{l+s+s1}{X1}\PY{l+s+s1}{\PYZsq{}}\PY{p}{]}
\PY{n}{X2} \PY{o}{=} \PY{n}{df}\PY{p}{[}\PY{l+s+s1}{\PYZsq{}}\PY{l+s+s1}{X2}\PY{l+s+s1}{\PYZsq{}}\PY{p}{]}
\PY{n}{X3} \PY{o}{=} \PY{n}{df}\PY{p}{[}\PY{l+s+s1}{\PYZsq{}}\PY{l+s+s1}{X3}\PY{l+s+s1}{\PYZsq{}}\PY{p}{]}
\PY{n}{X4} \PY{o}{=} \PY{n}{df}\PY{p}{[}\PY{l+s+s1}{\PYZsq{}}\PY{l+s+s1}{X4}\PY{l+s+s1}{\PYZsq{}}\PY{p}{]}

\PY{n+nb}{print}\PY{p}{(}\PY{l+s+sa}{f}\PY{l+s+s1}{\PYZsq{}}\PY{l+s+s1}{Минимум X1=}\PY{l+s+si}{\PYZob{}}\PY{n}{np}\PY{o}{.}\PY{n}{min}\PY{p}{(}\PY{n}{X1}\PY{p}{)}\PY{l+s+si}{\PYZcb{}}\PY{l+s+s1}{\PYZsq{}}\PY{p}{)}
\PY{n+nb}{print}\PY{p}{(}\PY{l+s+sa}{f}\PY{l+s+s1}{\PYZsq{}}\PY{l+s+s1}{Максимум X1=}\PY{l+s+si}{\PYZob{}}\PY{n}{np}\PY{o}{.}\PY{n}{max}\PY{p}{(}\PY{n}{X1}\PY{p}{)}\PY{l+s+si}{\PYZcb{}}\PY{l+s+s1}{\PYZsq{}}\PY{p}{)}
\PY{n+nb}{print}\PY{p}{(}\PY{p}{)}
\PY{n+nb}{print}\PY{p}{(}\PY{l+s+sa}{f}\PY{l+s+s1}{\PYZsq{}}\PY{l+s+s1}{Минимум X2=}\PY{l+s+si}{\PYZob{}}\PY{n}{np}\PY{o}{.}\PY{n}{min}\PY{p}{(}\PY{n}{X2}\PY{p}{)}\PY{l+s+si}{\PYZcb{}}\PY{l+s+s1}{\PYZsq{}}\PY{p}{)}
\PY{n+nb}{print}\PY{p}{(}\PY{l+s+sa}{f}\PY{l+s+s1}{\PYZsq{}}\PY{l+s+s1}{Максимум X2=}\PY{l+s+si}{\PYZob{}}\PY{n}{np}\PY{o}{.}\PY{n}{max}\PY{p}{(}\PY{n}{X2}\PY{p}{)}\PY{l+s+si}{\PYZcb{}}\PY{l+s+s1}{\PYZsq{}}\PY{p}{)}
\PY{n+nb}{print}\PY{p}{(}\PY{p}{)}
\PY{n+nb}{print}\PY{p}{(}\PY{l+s+sa}{f}\PY{l+s+s1}{\PYZsq{}}\PY{l+s+s1}{Минимум X3=}\PY{l+s+si}{\PYZob{}}\PY{n}{np}\PY{o}{.}\PY{n}{min}\PY{p}{(}\PY{n}{X3}\PY{p}{)}\PY{l+s+si}{\PYZcb{}}\PY{l+s+s1}{\PYZsq{}}\PY{p}{)}
\PY{n+nb}{print}\PY{p}{(}\PY{l+s+sa}{f}\PY{l+s+s1}{\PYZsq{}}\PY{l+s+s1}{Максимум X3=}\PY{l+s+si}{\PYZob{}}\PY{n}{np}\PY{o}{.}\PY{n}{max}\PY{p}{(}\PY{n}{X3}\PY{p}{)}\PY{l+s+si}{\PYZcb{}}\PY{l+s+s1}{\PYZsq{}}\PY{p}{)}
\PY{n+nb}{print}\PY{p}{(}\PY{p}{)}
\end{Verbatim}
\end{tcolorbox}

    \begin{Verbatim}[commandchars=\\\{\}]
Объем выборки n = 200
Минимум X1=28.81
Максимум X1=75.96

Минимум X2=26.8
Максимум X2=72.57

Минимум X3=22.47
Максимум X3=406.99

    \end{Verbatim}

    \begin{tcolorbox}[breakable, size=fbox, boxrule=1pt, pad at break*=1mm,colback=cellbackground, colframe=cellborder]
\prompt{In}{incolor}{5}{\boxspacing}
\begin{Verbatim}[commandchars=\\\{\}]
\PY{c+c1}{\PYZsh{} 4. Вариационный ряд \PYZhy{} просто сортировка}

\PY{k}{def}\PY{+w}{ }\PY{n+nf}{variation}\PY{p}{(}\PY{n}{x}\PY{p}{)}\PY{p}{:}
    \PY{n}{sorted\PYZus{}x} \PY{o}{=} \PY{n}{np}\PY{o}{.}\PY{n}{sort}\PY{p}{(}\PY{n}{x}\PY{p}{)}
    \PY{k}{return} \PY{n}{sorted\PYZus{}x}\PY{p}{[}\PY{p}{:}\PY{l+m+mi}{10}\PY{p}{]}  \PY{c+c1}{\PYZsh{} Вывести первые 10 для примера}


\PY{n+nb}{print}\PY{p}{(}\PY{n}{variation}\PY{p}{(}\PY{n}{X1}\PY{p}{)}\PY{p}{)}
\PY{n+nb}{print}\PY{p}{(}\PY{n}{variation}\PY{p}{(}\PY{n}{X2}\PY{p}{)}\PY{p}{)}
\PY{n+nb}{print}\PY{p}{(}\PY{n}{variation}\PY{p}{(}\PY{n}{X3}\PY{p}{)}\PY{p}{)}
\end{Verbatim}
\end{tcolorbox}

    \begin{Verbatim}[commandchars=\\\{\}]
[28.81 30.54 31.03 31.11 32.17 32.44 34.06 34.53 34.88 34.92]
[26.8  27.   27.53 27.57 27.62 28.39 28.61 28.71 29.35 29.58]
[22.47 22.76 22.89 22.98 23.53 23.8  24.05 24.47 25.21 25.68]
    \end{Verbatim}

    \begin{tcolorbox}[breakable, size=fbox, boxrule=1pt, pad at break*=1mm,colback=cellbackground, colframe=cellborder]
\prompt{In}{incolor}{6}{\boxspacing}
\begin{Verbatim}[commandchars=\\\{\}]
\PY{c+c1}{\PYZsh{}5. Эмпирическая функция распределения}

\PY{k}{def}\PY{+w}{ }\PY{n+nf}{empiric}\PY{p}{(}\PY{n}{x}\PY{p}{,} \PY{n}{place}\PY{p}{)}\PY{p}{:}
    \PY{n}{ecdf\PYZus{}x} \PY{o}{=} \PY{n}{ECDF}\PY{p}{(}\PY{n}{x}\PY{p}{)}
    \PY{n}{plt}\PY{o}{.}\PY{n}{step}\PY{p}{(}\PY{n}{ecdf\PYZus{}x}\PY{o}{.}\PY{n}{x}\PY{p}{,} \PY{n}{ecdf\PYZus{}x}\PY{o}{.}\PY{n}{y}\PY{p}{)}
    \PY{n}{plt}\PY{o}{.}\PY{n}{title}\PY{p}{(}\PY{l+s+sa}{f}\PY{l+s+s1}{\PYZsq{}}\PY{l+s+s1}{Эмпирическая функция распределения для }\PY{l+s+si}{\PYZob{}}\PY{n}{place}\PY{l+s+si}{\PYZcb{}}\PY{l+s+s1}{\PYZsq{}}\PY{p}{)}
    \PY{n}{plt}\PY{o}{.}\PY{n}{xlabel}\PY{p}{(}\PY{l+s+s1}{\PYZsq{}}\PY{l+s+s1}{x}\PY{l+s+s1}{\PYZsq{}}\PY{p}{)}
    \PY{n}{plt}\PY{o}{.}\PY{n}{ylabel}\PY{p}{(}\PY{l+s+s1}{\PYZsq{}}\PY{l+s+s1}{Fn(x)}\PY{l+s+s1}{\PYZsq{}}\PY{p}{)}
    \PY{n}{plt}\PY{o}{.}\PY{n}{grid}\PY{p}{(}\PY{p}{)}
    \PY{n}{plt}\PY{o}{.}\PY{n}{show}\PY{p}{(}\PY{p}{)}


\PY{n}{empiric}\PY{p}{(}\PY{n}{X1}\PY{p}{,} \PY{l+s+s1}{\PYZsq{}}\PY{l+s+s1}{X1}\PY{l+s+s1}{\PYZsq{}}\PY{p}{)}
\PY{n}{empiric}\PY{p}{(}\PY{n}{X2}\PY{p}{,} \PY{l+s+s1}{\PYZsq{}}\PY{l+s+s1}{X2}\PY{l+s+s1}{\PYZsq{}}\PY{p}{)}
\PY{n}{empiric}\PY{p}{(}\PY{n}{X3}\PY{p}{,} \PY{l+s+s1}{\PYZsq{}}\PY{l+s+s1}{X3}\PY{l+s+s1}{\PYZsq{}}\PY{p}{)}
\end{Verbatim}
\end{tcolorbox}

    \begin{center}
    \adjustimage{max size={0.9\linewidth}{0.9\paperheight}}{RGR1_A-2_files/RGR1_A-2_6_0.png}
    \end{center}
    { \hspace*{\fill} \\}
    
    \begin{center}
    \adjustimage{max size={0.9\linewidth}{0.9\paperheight}}{RGR1_A-2_files/RGR1_A-2_6_1.png}
    \end{center}
    { \hspace*{\fill} \\}
    
    \begin{center}
    \adjustimage{max size={0.9\linewidth}{0.9\paperheight}}{RGR1_A-2_files/RGR1_A-2_6_2.png}
    \end{center}
    { \hspace*{\fill} \\}
    
    \begin{tcolorbox}[breakable, size=fbox, boxrule=1pt, pad at break*=1mm,colback=cellbackground, colframe=cellborder]
\prompt{In}{incolor}{7}{\boxspacing}
\begin{Verbatim}[commandchars=\\\{\}]
\PY{c+c1}{\PYZsh{} 6. Гистограммы с разными правилами выбора интервалов}

\PY{k}{def}\PY{+w}{ }\PY{n+nf}{sterges}\PY{p}{(}\PY{n}{x}\PY{p}{,} \PY{n}{place}\PY{p}{)}\PY{p}{:}
    \PY{n}{n\PYZus{}bins\PYZus{}sturges} \PY{o}{=} \PY{n+nb}{int}\PY{p}{(}\PY{l+m+mi}{1} \PY{o}{+} \PY{n}{np}\PY{o}{.}\PY{n}{log2}\PY{p}{(}\PY{n}{n}\PY{p}{)}\PY{p}{)}
    \PY{n}{plt}\PY{o}{.}\PY{n}{hist}\PY{p}{(}\PY{n}{x}\PY{p}{,} \PY{n}{bins}\PY{o}{=}\PY{n}{n\PYZus{}bins\PYZus{}sturges}\PY{p}{,} \PY{n}{density}\PY{o}{=}\PY{k+kc}{True}\PY{p}{,} \PY{n}{alpha}\PY{o}{=}\PY{l+m+mf}{0.6}\PY{p}{,} \PY{n}{edgecolor}\PY{o}{=}\PY{l+s+s1}{\PYZsq{}}\PY{l+s+s1}{black}\PY{l+s+s1}{\PYZsq{}}\PY{p}{)}
    \PY{n}{plt}\PY{o}{.}\PY{n}{title}\PY{p}{(}\PY{l+s+sa}{f}\PY{l+s+s1}{\PYZsq{}}\PY{l+s+s1}{Гистограмма }\PY{l+s+si}{\PYZob{}}\PY{n}{place}\PY{l+s+si}{\PYZcb{}}\PY{l+s+s1}{ (Стерджес, bins=}\PY{l+s+si}{\PYZob{}}\PY{n}{n\PYZus{}bins\PYZus{}sturges}\PY{l+s+si}{\PYZcb{}}\PY{l+s+s1}{)}\PY{l+s+s1}{\PYZsq{}}\PY{p}{)}
    \PY{n}{plt}\PY{o}{.}\PY{n}{xlabel}\PY{p}{(}\PY{l+s+s1}{\PYZsq{}}\PY{l+s+s1}{x}\PY{l+s+s1}{\PYZsq{}}\PY{p}{)}
    \PY{n}{plt}\PY{o}{.}\PY{n}{ylabel}\PY{p}{(}\PY{l+s+s1}{\PYZsq{}}\PY{l+s+s1}{Плотность}\PY{l+s+s1}{\PYZsq{}}\PY{p}{)}
    \PY{n}{plt}\PY{o}{.}\PY{n}{grid}\PY{p}{(}\PY{k+kc}{True}\PY{p}{,} \PY{n}{alpha}\PY{o}{=}\PY{l+m+mf}{0.3}\PY{p}{)}
    \PY{n}{plt}\PY{o}{.}\PY{n}{show}\PY{p}{(}\PY{p}{)}


\PY{k}{def}\PY{+w}{ }\PY{n+nf}{scott}\PY{p}{(}\PY{n}{x}\PY{p}{,} \PY{n}{place}\PY{p}{)}\PY{p}{:}
    \PY{n}{h\PYZus{}scott} \PY{o}{=} \PY{l+m+mf}{3.5} \PY{o}{*} \PY{n}{np}\PY{o}{.}\PY{n}{std}\PY{p}{(}\PY{n}{x}\PY{p}{,} \PY{n}{ddof}\PY{o}{=}\PY{l+m+mi}{1}\PY{p}{)} \PY{o}{*} \PY{n}{n} \PY{o}{*}\PY{o}{*} \PY{p}{(}\PY{o}{\PYZhy{}}\PY{l+m+mi}{1} \PY{o}{/} \PY{l+m+mi}{3}\PY{p}{)}
    \PY{n}{n\PYZus{}bins\PYZus{}scott} \PY{o}{=} \PY{n+nb}{int}\PY{p}{(}\PY{n}{np}\PY{o}{.}\PY{n}{ceil}\PY{p}{(}\PY{p}{(}\PY{n}{x}\PY{o}{.}\PY{n}{max}\PY{p}{(}\PY{p}{)} \PY{o}{\PYZhy{}} \PY{n}{x}\PY{o}{.}\PY{n}{min}\PY{p}{(}\PY{p}{)}\PY{p}{)} \PY{o}{/} \PY{n}{h\PYZus{}scott}\PY{p}{)}\PY{p}{)}
    \PY{n}{plt}\PY{o}{.}\PY{n}{hist}\PY{p}{(}\PY{n}{x}\PY{p}{,} \PY{n}{bins}\PY{o}{=}\PY{n}{n\PYZus{}bins\PYZus{}scott}\PY{p}{,} \PY{n}{density}\PY{o}{=}\PY{k+kc}{True}\PY{p}{,} \PY{n}{alpha}\PY{o}{=}\PY{l+m+mf}{0.6}\PY{p}{,} \PY{n}{edgecolor}\PY{o}{=}\PY{l+s+s1}{\PYZsq{}}\PY{l+s+s1}{black}\PY{l+s+s1}{\PYZsq{}}\PY{p}{,} \PY{n}{color}\PY{o}{=}\PY{l+s+s1}{\PYZsq{}}\PY{l+s+s1}{orange}\PY{l+s+s1}{\PYZsq{}}\PY{p}{)}
    \PY{n}{plt}\PY{o}{.}\PY{n}{title}\PY{p}{(}\PY{l+s+sa}{f}\PY{l+s+s1}{\PYZsq{}}\PY{l+s+s1}{Гистограмма }\PY{l+s+si}{\PYZob{}}\PY{n}{place}\PY{l+s+si}{\PYZcb{}}\PY{l+s+s1}{ (Скотт, bins=}\PY{l+s+si}{\PYZob{}}\PY{n}{n\PYZus{}bins\PYZus{}scott}\PY{l+s+si}{\PYZcb{}}\PY{l+s+s1}{)}\PY{l+s+s1}{\PYZsq{}}\PY{p}{)}
    \PY{n}{plt}\PY{o}{.}\PY{n}{xlabel}\PY{p}{(}\PY{l+s+s1}{\PYZsq{}}\PY{l+s+s1}{x}\PY{l+s+s1}{\PYZsq{}}\PY{p}{)}
    \PY{n}{plt}\PY{o}{.}\PY{n}{ylabel}\PY{p}{(}\PY{l+s+s1}{\PYZsq{}}\PY{l+s+s1}{Плотность}\PY{l+s+s1}{\PYZsq{}}\PY{p}{)}
    \PY{n}{plt}\PY{o}{.}\PY{n}{grid}\PY{p}{(}\PY{k+kc}{True}\PY{p}{,} \PY{n}{alpha}\PY{o}{=}\PY{l+m+mf}{0.3}\PY{p}{)}
    \PY{n}{plt}\PY{o}{.}\PY{n}{show}\PY{p}{(}\PY{p}{)}


\PY{k}{def}\PY{+w}{ }\PY{n+nf}{fridman\PYZus{}diaconis}\PY{p}{(}\PY{n}{x}\PY{p}{,} \PY{n}{place}\PY{p}{)}\PY{p}{:}
    \PY{n}{iqr} \PY{o}{=} \PY{n}{np}\PY{o}{.}\PY{n}{percentile}\PY{p}{(}\PY{n}{x}\PY{p}{,} \PY{l+m+mi}{75}\PY{p}{)} \PY{o}{\PYZhy{}} \PY{n}{np}\PY{o}{.}\PY{n}{percentile}\PY{p}{(}\PY{n}{x}\PY{p}{,} \PY{l+m+mi}{25}\PY{p}{)}
    \PY{n}{h\PYZus{}fd} \PY{o}{=} \PY{l+m+mi}{2} \PY{o}{*} \PY{n}{iqr} \PY{o}{*} \PY{n}{n} \PY{o}{*}\PY{o}{*} \PY{p}{(}\PY{o}{\PYZhy{}}\PY{l+m+mi}{1} \PY{o}{/} \PY{l+m+mi}{3}\PY{p}{)}
    \PY{n}{n\PYZus{}bins\PYZus{}fd} \PY{o}{=} \PY{n+nb}{int}\PY{p}{(}\PY{n}{np}\PY{o}{.}\PY{n}{ceil}\PY{p}{(}\PY{p}{(}\PY{n}{x}\PY{o}{.}\PY{n}{max}\PY{p}{(}\PY{p}{)} \PY{o}{\PYZhy{}} \PY{n}{x}\PY{o}{.}\PY{n}{min}\PY{p}{(}\PY{p}{)}\PY{p}{)} \PY{o}{/} \PY{n}{h\PYZus{}fd}\PY{p}{)}\PY{p}{)}
    \PY{n}{plt}\PY{o}{.}\PY{n}{hist}\PY{p}{(}\PY{n}{x}\PY{p}{,} \PY{n}{bins}\PY{o}{=}\PY{n}{n\PYZus{}bins\PYZus{}fd}\PY{p}{,} \PY{n}{density}\PY{o}{=}\PY{k+kc}{True}\PY{p}{,} \PY{n}{alpha}\PY{o}{=}\PY{l+m+mf}{0.6}\PY{p}{,} \PY{n}{edgecolor}\PY{o}{=}\PY{l+s+s1}{\PYZsq{}}\PY{l+s+s1}{black}\PY{l+s+s1}{\PYZsq{}}\PY{p}{,} \PY{n}{color}\PY{o}{=}\PY{l+s+s1}{\PYZsq{}}\PY{l+s+s1}{green}\PY{l+s+s1}{\PYZsq{}}\PY{p}{)}
    \PY{n}{plt}\PY{o}{.}\PY{n}{title}\PY{p}{(}\PY{l+s+sa}{f}\PY{l+s+s1}{\PYZsq{}}\PY{l+s+s1}{Гистограмма }\PY{l+s+si}{\PYZob{}}\PY{n}{place}\PY{l+s+si}{\PYZcb{}}\PY{l+s+s1}{ (Фридман\PYZhy{}Диаконис, bins=}\PY{l+s+si}{\PYZob{}}\PY{n}{n\PYZus{}bins\PYZus{}fd}\PY{l+s+si}{\PYZcb{}}\PY{l+s+s1}{)}\PY{l+s+s1}{\PYZsq{}}\PY{p}{)}
    \PY{n}{plt}\PY{o}{.}\PY{n}{xlabel}\PY{p}{(}\PY{l+s+s1}{\PYZsq{}}\PY{l+s+s1}{x}\PY{l+s+s1}{\PYZsq{}}\PY{p}{)}
    \PY{n}{plt}\PY{o}{.}\PY{n}{ylabel}\PY{p}{(}\PY{l+s+s1}{\PYZsq{}}\PY{l+s+s1}{Плотность}\PY{l+s+s1}{\PYZsq{}}\PY{p}{)}
    \PY{n}{plt}\PY{o}{.}\PY{n}{grid}\PY{p}{(}\PY{k+kc}{True}\PY{p}{,} \PY{n}{alpha}\PY{o}{=}\PY{l+m+mf}{0.3}\PY{p}{)}
    \PY{n}{plt}\PY{o}{.}\PY{n}{show}\PY{p}{(}\PY{p}{)}


\PY{n}{sterges}\PY{p}{(}\PY{n}{X1}\PY{p}{,} \PY{l+s+s1}{\PYZsq{}}\PY{l+s+s1}{X1}\PY{l+s+s1}{\PYZsq{}}\PY{p}{)}
\PY{n}{scott}\PY{p}{(}\PY{n}{X1}\PY{p}{,} \PY{l+s+s1}{\PYZsq{}}\PY{l+s+s1}{X1}\PY{l+s+s1}{\PYZsq{}}\PY{p}{)}
\PY{n}{fridman\PYZus{}diaconis}\PY{p}{(}\PY{n}{X1}\PY{p}{,} \PY{l+s+s1}{\PYZsq{}}\PY{l+s+s1}{X1}\PY{l+s+s1}{\PYZsq{}}\PY{p}{)}

\PY{n}{sterges}\PY{p}{(}\PY{n}{X2}\PY{p}{,} \PY{l+s+s1}{\PYZsq{}}\PY{l+s+s1}{X2}\PY{l+s+s1}{\PYZsq{}}\PY{p}{)}
\PY{n}{scott}\PY{p}{(}\PY{n}{X2}\PY{p}{,} \PY{l+s+s1}{\PYZsq{}}\PY{l+s+s1}{X2}\PY{l+s+s1}{\PYZsq{}}\PY{p}{)}
\PY{n}{fridman\PYZus{}diaconis}\PY{p}{(}\PY{n}{X2}\PY{p}{,} \PY{l+s+s1}{\PYZsq{}}\PY{l+s+s1}{X2}\PY{l+s+s1}{\PYZsq{}}\PY{p}{)}

\PY{n}{sterges}\PY{p}{(}\PY{n}{X3}\PY{p}{,} \PY{l+s+s1}{\PYZsq{}}\PY{l+s+s1}{X3}\PY{l+s+s1}{\PYZsq{}}\PY{p}{)}
\PY{n}{scott}\PY{p}{(}\PY{n}{X3}\PY{p}{,} \PY{l+s+s1}{\PYZsq{}}\PY{l+s+s1}{X3}\PY{l+s+s1}{\PYZsq{}}\PY{p}{)}
\PY{n}{fridman\PYZus{}diaconis}\PY{p}{(}\PY{n}{X3}\PY{p}{,} \PY{l+s+s1}{\PYZsq{}}\PY{l+s+s1}{X3}\PY{l+s+s1}{\PYZsq{}}\PY{p}{)}
\end{Verbatim}
\end{tcolorbox}

    \begin{center}
    \adjustimage{max size={0.9\linewidth}{0.9\paperheight}}{RGR1_A-2_files/RGR1_A-2_7_0.png}
    \end{center}
    { \hspace*{\fill} \\}
    
    \begin{center}
    \adjustimage{max size={0.9\linewidth}{0.9\paperheight}}{RGR1_A-2_files/RGR1_A-2_7_1.png}
    \end{center}
    { \hspace*{\fill} \\}
    
    \begin{center}
    \adjustimage{max size={0.9\linewidth}{0.9\paperheight}}{RGR1_A-2_files/RGR1_A-2_7_2.png}
    \end{center}
    { \hspace*{\fill} \\}
    
    \begin{center}
    \adjustimage{max size={0.9\linewidth}{0.9\paperheight}}{RGR1_A-2_files/RGR1_A-2_7_3.png}
    \end{center}
    { \hspace*{\fill} \\}
    
    \begin{center}
    \adjustimage{max size={0.9\linewidth}{0.9\paperheight}}{RGR1_A-2_files/RGR1_A-2_7_4.png}
    \end{center}
    { \hspace*{\fill} \\}
    
    \begin{center}
    \adjustimage{max size={0.9\linewidth}{0.9\paperheight}}{RGR1_A-2_files/RGR1_A-2_7_5.png}
    \end{center}
    { \hspace*{\fill} \\}
    
    \begin{center}
    \adjustimage{max size={0.9\linewidth}{0.9\paperheight}}{RGR1_A-2_files/RGR1_A-2_7_6.png}
    \end{center}
    { \hspace*{\fill} \\}
    
    \begin{center}
    \adjustimage{max size={0.9\linewidth}{0.9\paperheight}}{RGR1_A-2_files/RGR1_A-2_7_7.png}
    \end{center}
    { \hspace*{\fill} \\}
    
    \begin{center}
    \adjustimage{max size={0.9\linewidth}{0.9\paperheight}}{RGR1_A-2_files/RGR1_A-2_7_8.png}
    \end{center}
    { \hspace*{\fill} \\}
    
    \begin{tcolorbox}[breakable, size=fbox, boxrule=1pt, pad at break*=1mm,colback=cellbackground, colframe=cellborder]
\prompt{In}{incolor}{8}{\boxspacing}
\begin{Verbatim}[commandchars=\\\{\}]
\PY{k}{def}\PY{+w}{ }\PY{n+nf}{numerical\PYZus{}characteristics}\PY{p}{(}\PY{n}{x}\PY{p}{,} \PY{n}{place}\PY{p}{)}\PY{p}{:}
    \PY{n}{x\PYZus{}bar} \PY{o}{=} \PY{n}{np}\PY{o}{.}\PY{n}{mean}\PY{p}{(}\PY{n}{x}\PY{p}{)}
    \PY{n}{S2} \PY{o}{=} \PY{n}{np}\PY{o}{.}\PY{n}{var}\PY{p}{(}\PY{n}{x}\PY{p}{,} \PY{n}{ddof}\PY{o}{=}\PY{l+m+mi}{0}\PY{p}{)}  \PY{c+c1}{\PYZsh{} по умолчанию смещенная}
    \PY{n}{sigma2\PYZus{}hat} \PY{o}{=} \PY{n}{np}\PY{o}{.}\PY{n}{var}\PY{p}{(}\PY{n}{x}\PY{p}{,} \PY{n}{ddof}\PY{o}{=}\PY{l+m+mi}{1}\PY{p}{)}  \PY{c+c1}{\PYZsh{} указатель на несмещенную}
    \PY{n}{S} \PY{o}{=} \PY{n}{np}\PY{o}{.}\PY{n}{sqrt}\PY{p}{(}\PY{n}{S2}\PY{p}{)}
    \PY{n}{sigma\PYZus{}hat} \PY{o}{=} \PY{n}{np}\PY{o}{.}\PY{n}{sqrt}\PY{p}{(}\PY{n}{sigma2\PYZus{}hat}\PY{p}{)}
    \PY{n}{med} \PY{o}{=} \PY{n}{np}\PY{o}{.}\PY{n}{median}\PY{p}{(}\PY{n}{x}\PY{p}{)}
    \PY{n}{q1} \PY{o}{=} \PY{n}{np}\PY{o}{.}\PY{n}{percentile}\PY{p}{(}\PY{n}{x}\PY{p}{,} \PY{l+m+mi}{25}\PY{p}{)}
    \PY{n}{q3} \PY{o}{=} \PY{n}{np}\PY{o}{.}\PY{n}{percentile}\PY{p}{(}\PY{n}{x}\PY{p}{,} \PY{l+m+mi}{75}\PY{p}{)}

    \PY{k}{return} \PY{l+s+sa}{f}\PY{l+s+s2}{\PYZdq{}}\PY{l+s+si}{\PYZob{}}\PY{n}{place}\PY{l+s+si}{\PYZcb{}}\PY{l+s+s2}{:}\PY{l+s+se}{\PYZbs{}n}\PY{l+s+s2}{ среднее=}\PY{l+s+si}{\PYZob{}}\PY{n}{x\PYZus{}bar}\PY{l+s+si}{:}\PY{l+s+s2}{.2f}\PY{l+s+si}{\PYZcb{}}\PY{l+s+se}{\PYZbs{}n}\PY{l+s+s2}{ медиана=}\PY{l+s+si}{\PYZob{}}\PY{n}{med}\PY{l+s+si}{:}\PY{l+s+s2}{.2f}\PY{l+s+si}{\PYZcb{}}\PY{l+s+se}{\PYZbs{}n}\PY{l+s+s2}{ смещенная дисперсия=}\PY{l+s+si}{\PYZob{}}\PY{n}{S2}\PY{l+s+si}{:}\PY{l+s+s2}{.2f}\PY{l+s+si}{\PYZcb{}}\PY{l+s+se}{\PYZbs{}n}\PY{l+s+s2}{ несмещенная дисперсия=}\PY{l+s+si}{\PYZob{}}\PY{n}{sigma2\PYZus{}hat}\PY{l+s+si}{:}\PY{l+s+s2}{.2f}\PY{l+s+si}{\PYZcb{}}\PY{l+s+se}{\PYZbs{}n}\PY{l+s+s2}{ смещенное выборочное стандартное отклонение=}\PY{l+s+si}{\PYZob{}}\PY{n}{S}\PY{l+s+si}{:}\PY{l+s+s2}{.2f}\PY{l+s+si}{\PYZcb{}}\PY{l+s+se}{\PYZbs{}n}\PY{l+s+s2}{ несмещенное выборочное стандартное отклонение=}\PY{l+s+si}{\PYZob{}}\PY{n}{sigma\PYZus{}hat}\PY{l+s+si}{:}\PY{l+s+s2}{.2f}\PY{l+s+si}{\PYZcb{}}\PY{l+s+se}{\PYZbs{}n}\PY{l+s+s2}{ медиана=}\PY{l+s+si}{\PYZob{}}\PY{n}{med}\PY{l+s+si}{:}\PY{l+s+s2}{.2f}\PY{l+s+si}{\PYZcb{}}\PY{l+s+se}{\PYZbs{}n}\PY{l+s+s2}{ первый квартиль \PYZhy{} 25\PYZhy{}й процентиль=}\PY{l+s+si}{\PYZob{}}\PY{n}{q1}\PY{l+s+si}{:}\PY{l+s+s2}{.2f}\PY{l+s+si}{\PYZcb{}}\PY{l+s+s2}{ \PYZhy{} 25\PYZpc{} данных меньше этого значения}\PY{l+s+se}{\PYZbs{}n}\PY{l+s+s2}{ третий квартиль \PYZhy{} 75\PYZhy{}й процентиль=}\PY{l+s+si}{\PYZob{}}\PY{n}{q3}\PY{l+s+si}{:}\PY{l+s+s2}{.2f}\PY{l+s+si}{\PYZcb{}}\PY{l+s+s2}{ \PYZhy{} 75\PYZpc{} данных меньше этого значения}\PY{l+s+se}{\PYZbs{}n}\PY{l+s+s2}{\PYZdq{}}


\PY{n+nb}{print}\PY{p}{(}\PY{n}{numerical\PYZus{}characteristics}\PY{p}{(}\PY{n}{X1}\PY{p}{,} \PY{l+s+s1}{\PYZsq{}}\PY{l+s+s1}{X1}\PY{l+s+s1}{\PYZsq{}}\PY{p}{)}\PY{p}{)}
\PY{n+nb}{print}\PY{p}{(}\PY{n}{numerical\PYZus{}characteristics}\PY{p}{(}\PY{n}{X2}\PY{p}{,} \PY{l+s+s1}{\PYZsq{}}\PY{l+s+s1}{X2}\PY{l+s+s1}{\PYZsq{}}\PY{p}{)}\PY{p}{)}
\PY{n+nb}{print}\PY{p}{(}\PY{n}{numerical\PYZus{}characteristics}\PY{p}{(}\PY{n}{X3}\PY{p}{,} \PY{l+s+s1}{\PYZsq{}}\PY{l+s+s1}{X3}\PY{l+s+s1}{\PYZsq{}}\PY{p}{)}\PY{p}{)}
\end{Verbatim}
\end{tcolorbox}

    \begin{Verbatim}[commandchars=\\\{\}]
X1:
 среднее=50.57
 медиана=50.34
 смещенная дисперсия=86.34
 несмещенная дисперсия=86.77
 смещенное выборочное стандартное отклонение=9.29
 несмещенное выборочное стандартное отклонение=9.32
 медиана=50.34
 первый квартиль - 25-й процентиль=44.39 - 25\% данных меньше этого значения
 третий квартиль - 75-й процентиль=56.56 - 75\% данных меньше этого значения

X2:
 среднее=49.10
 медиана=47.90
 смещенная дисперсия=166.74
 несмещенная дисперсия=167.58
 смещенное выборочное стандартное отклонение=12.91
 несмещенное выборочное стандартное отклонение=12.95
 медиана=47.90
 первый квартиль - 25-й процентиль=38.29 - 25\% данных меньше этого значения
 третий квартиль - 75-й процентиль=59.17 - 75\% данных меньше этого значения

X3:
 среднее=79.81
 медиана=63.35
 смещенная дисперсия=3301.58
 несмещенная дисперсия=3318.17
 смещенное выборочное стандартное отклонение=57.46
 несмещенное выборочное стандартное отклонение=57.60
 медиана=63.35
 первый квартиль - 25-й процентиль=41.16 - 25\% данных меньше этого значения
 третий квартиль - 75-й процентиль=98.18 - 75\% данных меньше этого значения

    \end{Verbatim}

    \paragraph{Описание формы
распределения}\label{ux43eux43fux438ux441ux430ux43dux438ux435-ux444ux43eux440ux43cux44b-ux440ux430ux441ux43fux440ux435ux434ux435ux43bux435ux43dux438ux44f}

\begin{itemize}
\tightlist
\item
  X1:

  \begin{itemize}
  \tightlist
  \item
    Распределение близко к симметричному. Среднее значение (50.57)
    практически совпадает с медианой (50.34), разница составляет всего
    0.23, что значительно меньше стандартного отклонения (9.32).
    Коэффициент асимметрии мал (0.27).
  \item
    Сильных выбросов практически нет. Межквартильный размах IQR = 12.17.
    По правилу \(Q_1 - 1.5 \cdot IQR;\ Q_3 + 1.5 \cdot IQR\) верхняя
    граница равна 74.81, поэтому только несколько верхних значений
    слегка выходят за нее; в целом основной массив наблюдений
    компактный.
  \item
    Явные естественные границы отсутствуют. Значения варьируются от
    28.81 до 75.96. Гистограмма не показывает жесткого обрыва ни слева,
    ни справа.
  \end{itemize}
\item
  X2:

  \begin{itemize}
  \tightlist
  \item
    Распределение близко к симметричному. Среднее значение (49.10)
    близко к медиане (47.90), разница составляет 1.20. Коэффициент
    асимметрии также мал (0.15).
  \item
    Выбросы отсутствуют. Межквартильный размах IQR = 20.89. Все значения
    выборки (min=26.80, max=72.57) находятся в интервале
    \([6.96;\ 90.51]\).
  \item
    Наблюдаются естественные нижняя и верхняя границы. По гистограмме
    видно, что данные заполняют конечный интервал примерно от 26.80 до
    72.57 без длинных хвостов.
  \end{itemize}
\item
  X3:

  \begin{itemize}
  \tightlist
  \item
    Наблюдается выраженная правосторонняя асимметрия. Среднее значение
    (79.81) заметно превышает медиану (63.35), разница составляет 16.46,
    что указывает на длинный правый хвост распределения.
  \item
    Выбросы присутствуют. Межквартильный размах IQR = 57.02. Верхняя
    граница по правилу \(Q_3 + 1.5 \cdot IQR\) равна 183.72, однако
    максимальное значение выборки равно 406.99, что значительно
    превышает эту границу.
  \item
    Имеется естественная нижняя граница. Минимальное значение выборки
    равно 22.47, все наблюдения лежат правее этой точки, при этом справа
    распределение имеет тяжелый хвост.
  \end{itemize}
\end{itemize}

    \subsubsection{2. Предположение о виде закона
распределения}\label{ux43fux440ux435ux434ux43fux43eux43bux43eux436ux435ux43dux438ux435-ux43e-ux432ux438ux434ux435-ux437ux430ux43aux43eux43dux430-ux440ux430ux441ux43fux440ux435ux434ux435ux43bux435ux43dux438ux44f}

    X1:

Гистограмма X1 имеет колоколообразную и почти симметричную форму, что
характерно для \textbf{нормального распределения}. Среднее (50.57) и
медиана (50.34) практически совпадают, а эмпирическая функция
распределения растет плавно без выраженных изломов. Небольшие отклонения
в правом хвосте не меняют общей картины. Таким образом, наиболее
подходящая модель - \(N_{a,\sigma}\)

X2:

Гистограмма X2 по форме близка к прямоугольной, а значения заполняют
конечный интервал примерно равномерно, что характерно для
\textbf{равномерного распределения}. Среднее (49.10) и медиана (47.90)
близки, выраженной асимметрии нет. Эмпирическая функция распределения на
основном промежутке близка к линейной. Таким образом, наиболее
подходящая модель - \(U_{a, b}\)

X3:

Гистограмма X3 имеет убывающую форму с длинным правым хвостом, что
характерно для \textbf{экспоненциального распределения со сдвигом}.
Среднее (79.81) существенно больше медианы (63.35), что указывает на
выраженную правостороннюю асимметрию. Все значения выборки лежат правее
точки 22.47, что естественно интерпретировать как сдвиг. Таким образом,
наиболее подходящая модель - \(Exp_{\lambda,c}\)

    \subsubsection{3. Оценивание параметров: метод моментов и метод
максимального
правдоподобия}\label{ux43eux446ux435ux43dux438ux432ux430ux43dux438ux435-ux43fux430ux440ux430ux43cux435ux442ux440ux43eux432-ux43cux435ux442ux43eux434-ux43cux43eux43cux435ux43dux442ux43eux432-ux438-ux43cux435ux442ux43eux434-ux43cux430ux43aux441ux438ux43cux430ux43bux44cux43dux43eux433ux43e-ux43fux440ux430ux432ux434ux43eux43fux43eux434ux43eux431ux438ux44f}

    \paragraph{X1: Нормальное
распределение}\label{x1-ux43dux43eux440ux43cux430ux43bux44cux43dux43eux435-ux440ux430ux441ux43fux440ux435ux434ux435ux43bux435ux43dux438ux435}

\begin{enumerate}
\def\labelenumi{\arabic{enumi}.}
\tightlist
\item
  Метод моментов
\end{enumerate}

Для нормального распределения теоретические моменты имеют вид:

\(EX = a, \ DX = \sigma^2\)

Приравнивая теоретические моменты к выборочным характеристикам

\(\hat{a}_{MM} = \bar{x} = 50.57\) - среднее арифметическое

\(\hat{\sigma}_{MM} = S = \sqrt{S^2} = \sqrt{86.34} = 9.29\) - смещенное
выборочное стандартное отклонение

\begin{enumerate}
\def\labelenumi{\arabic{enumi}.}
\setcounter{enumi}{1}
\tightlist
\item
  Метод максимального правдоподобия
\end{enumerate}

Функция правдоподобия:
\(L(a, \sigma) = \prod_{i = 1}^n \left(\frac{1}{\sigma\sqrt{2\pi}} \cdot \exp\left(-\frac{(x_i - a)^2}{2\sigma^2}\right)\right)\)

Упростим:
\(L(a, \sigma) = \left(\frac{1}{\sigma\sqrt{2\pi}}\right)^n \cdot \exp\left(-\frac{\sum (x_i - a)^2}{2\sigma^2}\right)\)

Логарифмическая функция правдоподобия:
\(l(a, \sigma) = \ln L(a, \sigma) = -n \ln(\sigma) - n \ln(\sqrt{2\pi}) - \frac{1}{2\sigma^2} \sum (x_i - a)^2\)

\textbf{Находим оценку} для \(a\) - берем частную производную по \(a\) и
приравниваем к 0

\(\frac{\partial l}{\partial a} = \frac{1}{\sigma^2} \sum (x_i - a) = 0\)

Значит
\(\sum (x_i - a) = 0 \Leftrightarrow a = \frac{\sum x_i}{n} = \bar{x}\)

\textbf{Находим оценку} для \(\sigma\) - берем частную производную по
\(\sigma\) и приравниваем к 0

\(\frac{\partial l}{\partial \sigma} = -\frac{n}{\sigma} + \frac{1}{\sigma^3} \sum (x_i - a)^2 = 0\)

Умножаем на
\(\sigma^3 \Rightarrow -n \sigma^2 + \sum (x_i - a)^2 = 0 \Rightarrow \sigma^2 = \frac{\sum (x_i - a)^2}{n} = S^2\)

Таким образом, получаем, что \(\sigma = \sqrt{S^2} = S\)

Соответственно, для нормального распределения оба метода дают идентичные
оценки. Это объясняется тем, что выборочное среднее и выборочная
дисперсия являются достаточными статистиками для параметров нормального
распределения.

    \paragraph{X2: Равномерное распределение U(a,
b)}\label{x2-ux440ux430ux432ux43dux43eux43cux435ux440ux43dux43eux435-ux440ux430ux441ux43fux440ux435ux434ux435ux43bux435ux43dux438ux435-ua-b}

Плотность распределения: \(f(x) = \frac{1}{b - a}, \ a \leq x \leq b\)

Теоретические моменты:
\(EX = \frac{a + b}{2}, \ DX = \frac{(b - a)^2}{12}\)

\begin{enumerate}
\def\labelenumi{\arabic{enumi}.}
\tightlist
\item
  Метод моментов
\end{enumerate}

Приравниваем теоретические моменты к выборочным характеристикам

\(\frac{a + b}{2} = \bar{x} = 49.10\)

\(\frac{(b - a)^2}{12} = S^2 = 166.74\)

Из второго уравнения получаем
\((b - a)^2 = 12 \cdot S^2 = 12 \cdot 166.74 = 2000.85 \Rightarrow b - a = \sqrt{12 \cdot S^2} = 44.73\)

При этом из первого уравнения получаем, что
\(a + b = 2 \bar{x} = 2 \cdot 49.10 = 98.20\)

Получается, решаем систему

\(\left\{\begin{array}{c}
a + b = 98.20 \\
b - a = 44.73
\end{array}\right. \Rightarrow 2b = 142.93 \Rightarrow b = 71.465 \Rightarrow a = 26.734\)

Таким образом, получаем, что
\(\hat{a}_{MM} = 26.734, \ \hat{b}_{MM} = 71.465\)

\begin{enumerate}
\def\labelenumi{\arabic{enumi}.}
\setcounter{enumi}{1}
\tightlist
\item
  Метод максимального правдоподобия
\end{enumerate}

Функция правдоподобия \(L(a, b) = \frac{1}{(b - a)^n}\) при условии
\(a \leq x_i \leq b\) для всех \(i\)

Для максимизации \(L(a, b)\) необходимо минимизировать \((b - a)\). Это
достигается при

\begin{itemize}
\tightlist
\item
  \(a = \max\) возможное, но не превышающее \(\min(x_i)\)
\item
  \(b = \min\) возможное, но не меньшее \(\max(x_i)\)
\end{itemize}

Следовательно,
\(\hat{a}_{MMP} = \min(x_i) = 26.80, \ \hat{b}_{MMP} = \max(x_i) = 72.57\)

Сравнение методов


\begin{longtable}[]{@{}llll@{}}
\toprule\noalign{}
Параметр & Метод моментов & ММП & Различие \\
\midrule\noalign{}
\endhead
\bottomrule\noalign{}
\endlastfoot
a & 26.734 & 26.80 & 0.066 \\
b & 71.465 & 72.57 & 1.105 \\
\end{longtable}
}

Оценки близки, но различаются. Метод моментов дает интервал
\((26.734; 71.465)\), а ММП - \((26.80; 72.57)\). Различие объясняется
тем, что ММП использует экстремальные статистики \(\min\) и \(\max\),
которые естественно определяют границы равномерного распределения.

    \paragraph{X3: Экспоненциальное распределение со сдвигом Exp(λ,
c)}\label{x3-ux44dux43aux441ux43fux43eux43dux435ux43dux446ux438ux430ux43bux44cux43dux43eux435-ux440ux430ux441ux43fux440ux435ux434ux435ux43bux435ux43dux438ux435-ux441ux43e-ux441ux434ux432ux438ux433ux43eux43c-expux3bb-c}

Плотность распределения -
\(f(x) = \lambda \cdot \exp\left(-\lambda(x - c)\right)\),
\(x \geq c, \ \lambda > 0\)

Теоретические моменты:

\begin{itemize}
\tightlist
\item
  \(EX = c + \frac{1}{\lambda}\)
\item
  \(DX = \frac{1}{\lambda^2}\)
\end{itemize}

\begin{enumerate}
\def\labelenumi{\arabic{enumi}.}
\tightlist
\item
  Метод моментов
\end{enumerate}

Приравниваем теоретические моменты к выборочным характеристикам

\(c + \frac{1}{\lambda} = \bar{x} = 79.81\)

\(\frac{1}{\lambda^2} = S^2 = 3301.58\)

Из второго уравнения получаем
\(\lambda = \frac{1}{\sqrt{S^2}} = \frac{1}{S} = \frac{1}{57.46} = 0.0174\)

Из первого уравнения получаем
\(c = \bar{x} - \frac{1}{\lambda} = 79.81 - 57.46 = 22.35\)

Таким образом получаем, что:

\(\hat{c}_{MM} = 22.35, \ \hat{\lambda}_{MM} = 0.0174\)

\begin{enumerate}
\def\labelenumi{\arabic{enumi}.}
\setcounter{enumi}{1}
\tightlist
\item
  Метод максимального правдоподобия
\end{enumerate}

\(L(\lambda, c) = \lambda^n \cdot \exp\left(-\lambda \cdot \sum_{i = 1}^n (x_i - c)\right)\)
при условии \(x_i \geq c\)

Логарифмическая функция правдоподобия

\(l(\lambda, c) = n \cdot \ln(\lambda) - \lambda \cdot \sum_{i = 1}^n (x_i - c)\)

Для максимизации \(l(\lambda, c)\) необходимо выбрать \(c\) как можно
больше, но при этом есть условие, что
\(c \leq \min(x_i) \Rightarrow \hat{c}_{MMP} = \min(x_i) = 22.47\)

При фиксированном \(\hat{c}\) оценка \(\lambda\) находится из условия
\(\frac{\partial l}{\partial \lambda} = \frac{n}{\lambda} - \sum(x_i - \hat{c}) = 0\)

Тогда
\(\hat{\lambda}_{MMP} = \frac{n}{\sum (x_i - \hat{c})} = \frac{1}{\bar{x} - \hat{c}} = \frac{1}{79.81 - 22.47} = 0.0174\)

Сравнение методов:

{\def\LTcaptype{none} % do not increment counter
\begin{longtable}[]{@{}llll@{}}
\toprule\noalign{}
Параметр & Метод моментов & ММП & Различие \\
\midrule\noalign{}
\endhead
\bottomrule\noalign{}
\endlastfoot
\(c\) & 22.35 & 22.47 & 0.12 \\
\(\lambda\) & 0.0174 & 0.0174 & 0.0000 \\
\end{longtable}
}

Оценки практически совпадают. Метод моментов дает оценку сдвига
\(\hat{c} = 22.35\), а ММП оценивает сдвиг как минимальное значение
выборки \(\hat{c} = 22.47\). Такое малое расхождение означает, что
выбранная экспоненциальная модель хорошо согласуется с данными.

    \subsubsection{4. Оценивание параметрической вероятности двумя
способами}\label{ux43eux446ux435ux43dux438ux432ux430ux43dux438ux435-ux43fux430ux440ux430ux43cux435ux442ux440ux438ux447ux435ux441ux43aux43eux439-ux432ux435ux440ux43eux44fux442ux43dux43eux441ux442ux438-ux434ux432ux443ux43cux44f-ux441ux43fux43eux441ux43eux431ux430ux43cux438}

    \paragraph{Формулы для каждого
распределения}\label{ux444ux43eux440ux43cux443ux43bux44b-ux434ux43bux44f-ux43aux430ux436ux434ux43eux433ux43e-ux440ux430ux441ux43fux440ux435ux434ux435ux43bux435ux43dux438ux44f}

{\def\LTcaptype{none} % do not increment counter
\begin{longtable}[]{@{}
  >{\raggedright\arraybackslash}p{(\linewidth - 2\tabcolsep) * \real{0.1791}}
  >{\raggedright\arraybackslash}p{(\linewidth - 2\tabcolsep) * \real{0.8209}}@{}}
\toprule\noalign{}
\begin{minipage}[b]{\linewidth}\raggedright
Распределение
\end{minipage} & \begin{minipage}[b]{\linewidth}\raggedright
\(P(X > x_0)\)
\end{minipage} \\
\midrule\noalign{}
\endhead
\bottomrule\noalign{}
\endlastfoot
Нормальное \(N(a, \sigma)\) &
\(1 - \Phi\left(\frac{x_0 - a}{\sigma}\right)\), где \(\Phi\) - функция
стандартного нормального распределения \\
Равномерное \(U(a, b)\) &
\(\left\{\begin{array}{c} \frac{b - x_0}{b - a}, \ \text{если} \ a \leq x_0 \leq b; \\ 1, \ \text{если} \ x_0 < a; \\ 0, \ \text{если} \ x_0 > b \end{array}\right.\) \\
Экспоненциальное \(Exp(\lambda, c)\) &
\(\left\{\begin{array}{c} e^{-\lambda(x_0 - c)}, \ \text{если} \ x_0 \geq c; \\ 1, \ \text{если} \ x_0 < c \end{array}\right.\) \\
\end{longtable}
}

    \paragraph{Расчет для X1 - нормального
распределения}\label{ux440ux430ux441ux447ux435ux442-ux434ux43bux44f-x1---ux43dux43eux440ux43cux430ux43bux44cux43dux43eux433ux43e-ux440ux430ux441ux43fux440ux435ux434ux435ux43bux435ux43dux438ux44f}

Исходные данные:

\begin{itemize}
\tightlist
\item
  \(\bar{x} = 50.57\)
\item
  \(\hat{\sigma} = 9.32\) - несмещенное выборочное стандартное
  отклонение
\end{itemize}

    \textbf{Выбираем порог}

\(x_0 = \bar{x} + \hat{\sigma} = 50.57 + 9.32 = 59.88\)

    \begin{tcolorbox}[breakable, size=fbox, boxrule=1pt, pad at break*=1mm,colback=cellbackground, colframe=cellborder]
\prompt{In}{incolor}{9}{\boxspacing}
\begin{Verbatim}[commandchars=\\\{\}]
\PY{c+c1}{\PYZsh{} Эмпирическая оценка}

\PY{n}{p\PYZus{}emp} \PY{o}{=} \PY{n+nb}{sum}\PY{p}{(}\PY{n}{X1} \PY{o}{\PYZgt{}} \PY{l+m+mf}{59.88395}\PY{p}{)} \PY{o}{/} \PY{n}{n}

\PY{n+nb}{print}\PY{p}{(}\PY{n}{p\PYZus{}emp}\PY{p}{)}
\end{Verbatim}
\end{tcolorbox}

    \begin{Verbatim}[commandchars=\\\{\}]
0.14
    \end{Verbatim}

    Параметрическая оценка

\(P(X > x_0) = 1 - \Phi\left(\frac{x_0 - a}{\sigma}\right) = 1 - \Phi\left(\frac{59.88 - 50.57}{9.29}\right) = 1 - \Phi\left(1.003\right) \approx 0.1580\)

    \begin{tcolorbox}[breakable, size=fbox, boxrule=1pt, pad at break*=1mm,colback=cellbackground, colframe=cellborder]
\prompt{In}{incolor}{10}{\boxspacing}
\begin{Verbatim}[commandchars=\\\{\}]
\PY{c+c1}{\PYZsh{} Вычисление функции распределения стандартного нормального закона}
\PY{k+kn}{import}\PY{+w}{ }\PY{n+nn}{scipy}\PY{n+nn}{.}\PY{n+nn}{stats}\PY{+w}{ }\PY{k}{as}\PY{+w}{ }\PY{n+nn}{stats}

\PY{n+nb}{print}\PY{p}{(}\PY{n}{stats}\PY{o}{.}\PY{n}{norm}\PY{o}{.}\PY{n}{cdf}\PY{p}{(}\PY{l+m+mi}{1}\PY{p}{)}\PY{p}{)}
\end{Verbatim}
\end{tcolorbox}

    \begin{Verbatim}[commandchars=\\\{\}]
0.8413447460685429
    \end{Verbatim}

    Таким образом, получаем, что расхождение между эмпирической (0.14) и
параметрической (0.1580) оценками составляет около 1.8\%, что
свидетельствует о хорошем согласии выбранной нормальной модели с данными

    \paragraph{Расчет для X2 - равномерного
распределения}\label{ux440ux430ux441ux447ux435ux442-ux434ux43bux44f-x2---ux440ux430ux432ux43dux43eux43cux435ux440ux43dux43eux433ux43e-ux440ux430ux441ux43fux440ux435ux434ux435ux43bux435ux43dux438ux44f}

Исходные данные:

\begin{itemize}
\tightlist
\item
  \(\bar{x} = 49.10\)
\item
  \(\hat{\sigma} = 12.95\) - несмещенное выборочное стандартное
  отклонение
\item
  \(\hat{a} = 26.80\) - метод максимального правдоподобия
\item
  \(\hat{b} = 72.57\) - ММП
\end{itemize}

    \textbf{Выбираем порог}

\(x_0 = \bar{x} + \hat{\sigma} = 49.10 + 12.95 = 62.04\)

    \begin{tcolorbox}[breakable, size=fbox, boxrule=1pt, pad at break*=1mm,colback=cellbackground, colframe=cellborder]
\prompt{In}{incolor}{11}{\boxspacing}
\begin{Verbatim}[commandchars=\\\{\}]
\PY{c+c1}{\PYZsh{} Эмпирическая оценка}

\PY{n}{p\PYZus{}emp} \PY{o}{=} \PY{n+nb}{sum}\PY{p}{(}\PY{n}{X2} \PY{o}{\PYZgt{}} \PY{l+m+mf}{62.04459}\PY{p}{)} \PY{o}{/} \PY{n}{n}

\PY{n+nb}{print}\PY{p}{(}\PY{n}{p\PYZus{}emp}\PY{p}{)}
\end{Verbatim}
\end{tcolorbox}

    \begin{Verbatim}[commandchars=\\\{\}]
0.205
    \end{Verbatim}

    Параметрическая оценка

\(P(X > x_0) = \frac{b - x_0}{b - a} = \frac{72.57 - 62.04}{72.57 - 26.80} = \frac{10.53}{45.77} = 0.230\)

    Расхождение между эмпирической (0.205) и параметрической (0.230)
оценками составляет около 2.5\%, что указывает на приемлемое согласие
равномерной модели с данными

    \subsubsection{Расчет для X3 - экспоненциального
распределения}\label{ux440ux430ux441ux447ux435ux442-ux434ux43bux44f-x3---ux44dux43aux441ux43fux43eux43dux435ux43dux446ux438ux430ux43bux44cux43dux43eux433ux43e-ux440ux430ux441ux43fux440ux435ux434ux435ux43bux435ux43dux438ux44f}

Исходные данные:

\begin{itemize}
\tightlist
\item
  \(\bar{x} = 79.81\)
\item
  \(\hat{\sigma} = 57.60\) - несмещенное выборочное стандартное
  отклонение
\item
  \(\hat{c} = 22.47\) - метод максимального правдоподобия
\item
  \(\hat{\lambda} = 0.0174\) - ММП
\end{itemize}

    \textbf{Выбираем порог}

\(x_0 = \bar{x} + \hat{\sigma} = 79.81 + 57.60 = 137.41\)

    \begin{tcolorbox}[breakable, size=fbox, boxrule=1pt, pad at break*=1mm,colback=cellbackground, colframe=cellborder]
\prompt{In}{incolor}{12}{\boxspacing}
\begin{Verbatim}[commandchars=\\\{\}]
\PY{c+c1}{\PYZsh{} Эмпирическая оценка}

\PY{n}{p\PYZus{}emp} \PY{o}{=} \PY{n+nb}{sum}\PY{p}{(}\PY{n}{X3} \PY{o}{\PYZgt{}} \PY{l+m+mf}{137.41242}\PY{p}{)} \PY{o}{/} \PY{n}{n}

\PY{n+nb}{print}\PY{p}{(}\PY{n}{p\PYZus{}emp}\PY{p}{)}
\end{Verbatim}
\end{tcolorbox}

    \begin{Verbatim}[commandchars=\\\{\}]
0.135
    \end{Verbatim}

    Параметрическая оценка

\(P(X > x_0) = \exp\left(-\lambda(x_0 - c)\right) = \exp\left(-0.0174 \cdot (137.41 - 22.47)\right) = 0.1347\)

    \begin{tcolorbox}[breakable, size=fbox, boxrule=1pt, pad at break*=1mm,colback=cellbackground, colframe=cellborder]
\prompt{In}{incolor}{13}{\boxspacing}
\begin{Verbatim}[commandchars=\\\{\}]
\PY{k+kn}{from}\PY{+w}{ }\PY{n+nn}{math}\PY{+w}{ }\PY{k+kn}{import} \PY{n}{e}

\PY{n+nb}{print}\PY{p}{(}\PY{n}{e} \PY{o}{*}\PY{o}{*} \PY{p}{(}\PY{o}{\PYZhy{}}\PY{l+m+mf}{0.0174401671} \PY{o}{*} \PY{p}{(}\PY{l+m+mf}{137.41242} \PY{o}{\PYZhy{}} \PY{l+m+mf}{22.47}\PY{p}{)}\PY{p}{)}\PY{p}{)}
\end{Verbatim}
\end{tcolorbox}

    \begin{Verbatim}[commandchars=\\\{\}]
0.13471214831742076
    \end{Verbatim}

    Эмпирическая и параметрическая оценки вероятности \(P(X > x_0)\)
практически совпали: 0.135 и 0.1347. Это свидетельствует о хорошем
согласии выбранной экспоненциальной модели с данными и подтверждает
адекватность полученных оценок параметров

    \subsubsection{5. Оценка моментов по сгруппированной
выборке}\label{ux43eux446ux435ux43dux43aux430-ux43cux43eux43cux435ux43dux442ux43eux432-ux43fux43e-ux441ux433ux440ux443ux43fux43fux438ux440ux43eux432ux430ux43dux43dux43eux439-ux432ux44bux431ux43eux440ux43aux435}

\(\hat{X}_g = \frac{1}{n}\sum_{k=1}^m n_k\hat{X}_k\)

\(\hat{\sigma}_g^2 = \frac{1}{n - 1}\sum_{k=1}^m n_k(\hat{X}_k - \hat{X}_g)^2\)

где

\begin{itemize}
\tightlist
\item
  \(n\) - объем выборки
\item
  \(n_k\) - частота в k-м интервале
\item
  \(\hat{X}_k\) - середина k-го интервала
\end{itemize}

    По правилу Стерджеса
\(k = \lceil 1 + \log_2(n) \rceil = \lceil 1 + \log_2(200) \rceil = \lceil 8.64 \rceil = 9\)

    \begin{tcolorbox}[breakable, size=fbox, boxrule=1pt, pad at break*=1mm,colback=cellbackground, colframe=cellborder]
\prompt{In}{incolor}{14}{\boxspacing}
\begin{Verbatim}[commandchars=\\\{\}]
\PY{k+kn}{from}\PY{+w}{ }\PY{n+nn}{math}\PY{+w}{ }\PY{k+kn}{import} \PY{n}{log2}\PY{p}{,} \PY{n}{ceil}

\PY{n}{k} \PY{o}{=} \PY{n}{ceil}\PY{p}{(}\PY{l+m+mi}{1} \PY{o}{+} \PY{n}{log2}\PY{p}{(}\PY{n}{n}\PY{p}{)}\PY{p}{)}

\PY{n+nb}{print}\PY{p}{(}\PY{n}{k}\PY{p}{)}
\end{Verbatim}
\end{tcolorbox}

    \begin{Verbatim}[commandchars=\\\{\}]
9
    \end{Verbatim}

    \begin{tcolorbox}[breakable, size=fbox, boxrule=1pt, pad at break*=1mm,colback=cellbackground, colframe=cellborder]
\prompt{In}{incolor}{15}{\boxspacing}
\begin{Verbatim}[commandchars=\\\{\}]
\PY{n}{n} \PY{o}{=} \PY{l+m+mi}{200}
\PY{n}{k} \PY{o}{=} \PY{l+m+mi}{9}

\PY{k}{def}\PY{+w}{ }\PY{n+nf}{moments}\PY{p}{(}\PY{n}{x}\PY{p}{,} \PY{n}{place}\PY{p}{)}\PY{p}{:}
    \PY{c+c1}{\PYZsh{} Построение гистограммы}
    \PY{n}{counts}\PY{p}{,} \PY{n}{bin\PYZus{}edges} \PY{o}{=} \PY{n}{np}\PY{o}{.}\PY{n}{histogram}\PY{p}{(}\PY{n}{x}\PY{p}{,} \PY{n}{bins}\PY{o}{=}\PY{n}{k}\PY{p}{)}

    \PY{c+c1}{\PYZsh{} Середины интервалов}
    \PY{n}{bin\PYZus{}centers} \PY{o}{=} \PY{p}{(}\PY{n}{bin\PYZus{}edges}\PY{p}{[}\PY{p}{:}\PY{o}{\PYZhy{}}\PY{l+m+mi}{1}\PY{p}{]} \PY{o}{+} \PY{n}{bin\PYZus{}edges}\PY{p}{[}\PY{l+m+mi}{1}\PY{p}{:}\PY{p}{]}\PY{p}{)} \PY{o}{/} \PY{l+m+mi}{2}

    \PY{n}{all\PYZus{}counts} \PY{o}{=} \PY{l+m+mi}{0}
    \PY{n}{all\PYZus{}product} \PY{o}{=} \PY{l+m+mi}{0}

    \PY{n+nb}{print}\PY{p}{(}\PY{l+s+sa}{f}\PY{l+s+s2}{\PYZdq{}}\PY{l+s+s2}{Интервалы для }\PY{l+s+si}{\PYZob{}}\PY{n}{place}\PY{l+s+si}{\PYZcb{}}\PY{l+s+s2}{:}\PY{l+s+s2}{\PYZdq{}}\PY{p}{)}
    \PY{k}{for} \PY{n}{i} \PY{o+ow}{in} \PY{n+nb}{range}\PY{p}{(}\PY{n}{k}\PY{p}{)}\PY{p}{:}
        \PY{n}{all\PYZus{}counts} \PY{o}{+}\PY{o}{=} \PY{n}{counts}\PY{p}{[}\PY{n}{i}\PY{p}{]}
        \PY{n}{prod} \PY{o}{=} \PY{n}{counts}\PY{p}{[}\PY{n}{i}\PY{p}{]} \PY{o}{*} \PY{n}{bin\PYZus{}centers}\PY{p}{[}\PY{n}{i}\PY{p}{]}
        \PY{n}{all\PYZus{}product} \PY{o}{+}\PY{o}{=} \PY{n}{prod}
        \PY{n+nb}{print}\PY{p}{(}
            \PY{l+s+sa}{f}\PY{l+s+s2}{\PYZdq{}}\PY{l+s+si}{\PYZob{}}\PY{n}{i}\PY{+w}{ }\PY{o}{+}\PY{+w}{ }\PY{l+m+mi}{1}\PY{l+s+si}{\PYZcb{}}\PY{l+s+s2}{: [}\PY{l+s+si}{\PYZob{}}\PY{n}{bin\PYZus{}edges}\PY{p}{[}\PY{n}{i}\PY{p}{]}\PY{l+s+si}{:}\PY{l+s+s2}{.2f}\PY{l+s+si}{\PYZcb{}}\PY{l+s+s2}{, }\PY{l+s+si}{\PYZob{}}\PY{n}{bin\PYZus{}edges}\PY{p}{[}\PY{n}{i}\PY{+w}{ }\PY{o}{+}\PY{+w}{ }\PY{l+m+mi}{1}\PY{p}{]}\PY{l+s+si}{:}\PY{l+s+s2}{.2f}\PY{l+s+si}{\PYZcb{}}\PY{l+s+s2}{] | }\PY{l+s+s2}{\PYZdq{}}
            \PY{l+s+sa}{f}\PY{l+s+s2}{\PYZdq{}}\PY{l+s+s2}{частота = }\PY{l+s+si}{\PYZob{}}\PY{n}{counts}\PY{p}{[}\PY{n}{i}\PY{p}{]}\PY{l+s+si}{\PYZcb{}}\PY{l+s+s2}{ | середина = }\PY{l+s+si}{\PYZob{}}\PY{n}{bin\PYZus{}centers}\PY{p}{[}\PY{n}{i}\PY{p}{]}\PY{l+s+si}{:}\PY{l+s+s2}{.2f}\PY{l+s+si}{\PYZcb{}}\PY{l+s+s2}{ | }\PY{l+s+s2}{\PYZdq{}}
            \PY{l+s+sa}{f}\PY{l+s+s2}{\PYZdq{}}\PY{l+s+s2}{n\PYZus{}k * X\PYZus{}k = }\PY{l+s+si}{\PYZob{}}\PY{n}{prod}\PY{l+s+si}{:}\PY{l+s+s2}{.2f}\PY{l+s+si}{\PYZcb{}}\PY{l+s+s2}{\PYZdq{}}
        \PY{p}{)}

    \PY{n+nb}{print}\PY{p}{(}\PY{l+s+sa}{f}\PY{l+s+s2}{\PYZdq{}}\PY{l+s+s2}{Итог сумм: частот = }\PY{l+s+si}{\PYZob{}}\PY{n}{all\PYZus{}counts}\PY{l+s+si}{\PYZcb{}}\PY{l+s+s2}{, n\PYZus{}k * X\PYZus{}k = }\PY{l+s+si}{\PYZob{}}\PY{n}{all\PYZus{}product}\PY{l+s+si}{:}\PY{l+s+s2}{.2f}\PY{l+s+si}{\PYZcb{}}\PY{l+s+s2}{\PYZdq{}}\PY{p}{)}
    \PY{n}{x\PYZus{}g} \PY{o}{=} \PY{n}{all\PYZus{}product} \PY{o}{/} \PY{n}{all\PYZus{}counts}
    \PY{n}{display}\PY{p}{(}\PY{n}{Math}\PY{p}{(}\PY{l+s+sa}{rf}\PY{l+s+s2}{\PYZdq{}}\PY{l+s+s2}{\PYZbs{}}\PY{l+s+s2}{hat}\PY{l+s+se}{\PYZob{}\PYZob{}}\PY{l+s+s2}{X}\PY{l+s+se}{\PYZcb{}\PYZcb{}}\PY{l+s+s2}{\PYZus{}g = }\PY{l+s+s2}{\PYZbs{}}\PY{l+s+s2}{frac}\PY{l+s+se}{\PYZob{}\PYZob{}}\PY{l+s+si}{\PYZob{}}\PY{n}{all\PYZus{}product}\PY{l+s+si}{:}\PY{l+s+s2}{.2f}\PY{l+s+se}{\PYZcb{}\PYZcb{}}\PY{l+s+si}{\PYZcb{}}\PY{l+s+se}{\PYZob{}\PYZob{}}\PY{l+s+si}{\PYZob{}}\PY{n}{all\PYZus{}counts}\PY{l+s+si}{\PYZcb{}}\PY{l+s+se}{\PYZcb{}\PYZcb{}}\PY{l+s+s2}{ = }\PY{l+s+si}{\PYZob{}}\PY{n}{x\PYZus{}g}\PY{l+s+si}{:}\PY{l+s+s2}{.2f}\PY{l+s+si}{\PYZcb{}}\PY{l+s+s2}{\PYZdq{}}\PY{p}{)}\PY{p}{)}

    \PY{n+nb}{print}\PY{p}{(}\PY{p}{)}

    \PY{c+c1}{\PYZsh{} Дисперсия по сгруппированным данным}
    \PY{n}{all\PYZus{}disp} \PY{o}{=} \PY{l+m+mi}{0}
    \PY{k}{for} \PY{n}{i} \PY{o+ow}{in} \PY{n+nb}{range}\PY{p}{(}\PY{n}{k}\PY{p}{)}\PY{p}{:}
        \PY{n}{prod} \PY{o}{=} \PY{n}{counts}\PY{p}{[}\PY{n}{i}\PY{p}{]} \PY{o}{*} \PY{p}{(}\PY{n}{bin\PYZus{}centers}\PY{p}{[}\PY{n}{i}\PY{p}{]} \PY{o}{\PYZhy{}} \PY{n}{x\PYZus{}g}\PY{p}{)} \PY{o}{*}\PY{o}{*} \PY{l+m+mi}{2}
        \PY{n}{all\PYZus{}disp} \PY{o}{+}\PY{o}{=} \PY{n}{prod}
        \PY{n+nb}{print}\PY{p}{(}
            \PY{l+s+sa}{f}\PY{l+s+s2}{\PYZdq{}}\PY{l+s+si}{\PYZob{}}\PY{n}{i}\PY{+w}{ }\PY{o}{+}\PY{+w}{ }\PY{l+m+mi}{1}\PY{l+s+si}{\PYZcb{}}\PY{l+s+s2}{: [}\PY{l+s+si}{\PYZob{}}\PY{n}{bin\PYZus{}edges}\PY{p}{[}\PY{n}{i}\PY{p}{]}\PY{l+s+si}{:}\PY{l+s+s2}{.2f}\PY{l+s+si}{\PYZcb{}}\PY{l+s+s2}{, }\PY{l+s+si}{\PYZob{}}\PY{n}{bin\PYZus{}edges}\PY{p}{[}\PY{n}{i}\PY{+w}{ }\PY{o}{+}\PY{+w}{ }\PY{l+m+mi}{1}\PY{p}{]}\PY{l+s+si}{:}\PY{l+s+s2}{.2f}\PY{l+s+si}{\PYZcb{}}\PY{l+s+s2}{] | }\PY{l+s+s2}{\PYZdq{}}
            \PY{l+s+sa}{f}\PY{l+s+s2}{\PYZdq{}}\PY{l+s+s2}{частота = }\PY{l+s+si}{\PYZob{}}\PY{n}{counts}\PY{p}{[}\PY{n}{i}\PY{p}{]}\PY{l+s+si}{\PYZcb{}}\PY{l+s+s2}{ | середина = }\PY{l+s+si}{\PYZob{}}\PY{n}{bin\PYZus{}centers}\PY{p}{[}\PY{n}{i}\PY{p}{]}\PY{l+s+si}{:}\PY{l+s+s2}{.2f}\PY{l+s+si}{\PYZcb{}}\PY{l+s+s2}{ | }\PY{l+s+s2}{\PYZdq{}}
            \PY{l+s+sa}{f}\PY{l+s+s2}{\PYZdq{}}\PY{l+s+s2}{n\PYZus{}k * (X\PYZus{}k \PYZhy{} X\PYZus{}g)\PYZca{}2 = }\PY{l+s+si}{\PYZob{}}\PY{n}{prod}\PY{l+s+si}{:}\PY{l+s+s2}{.2f}\PY{l+s+si}{\PYZcb{}}\PY{l+s+s2}{\PYZdq{}}
        \PY{p}{)}

    \PY{n}{s2\PYZus{}g} \PY{o}{=} \PY{n}{all\PYZus{}disp} \PY{o}{/} \PY{p}{(}\PY{n}{all\PYZus{}counts} \PY{o}{\PYZhy{}} \PY{l+m+mi}{1}\PY{p}{)}
    \PY{n}{display}\PY{p}{(}\PY{n}{Math}\PY{p}{(}\PY{l+s+sa}{rf}\PY{l+s+s2}{\PYZdq{}}\PY{l+s+s2}{\PYZbs{}}\PY{l+s+s2}{hat}\PY{l+s+se}{\PYZob{}\PYZob{}}\PY{l+s+s2}{\PYZbs{}}\PY{l+s+s2}{sigma}\PY{l+s+se}{\PYZcb{}\PYZcb{}}\PY{l+s+s2}{\PYZus{}g\PYZca{}2 = }\PY{l+s+s2}{\PYZbs{}}\PY{l+s+s2}{frac}\PY{l+s+se}{\PYZob{}\PYZob{}}\PY{l+s+si}{\PYZob{}}\PY{n}{all\PYZus{}disp}\PY{l+s+si}{:}\PY{l+s+s2}{.2f}\PY{l+s+se}{\PYZcb{}\PYZcb{}}\PY{l+s+si}{\PYZcb{}}\PY{l+s+se}{\PYZob{}\PYZob{}}\PY{l+s+si}{\PYZob{}}\PY{n}{all\PYZus{}counts}\PY{+w}{ }\PY{o}{\PYZhy{}}\PY{+w}{ }\PY{l+m+mi}{1}\PY{l+s+si}{\PYZcb{}}\PY{l+s+se}{\PYZcb{}\PYZcb{}}\PY{l+s+s2}{ = }\PY{l+s+si}{\PYZob{}}\PY{n}{s2\PYZus{}g}\PY{l+s+si}{:}\PY{l+s+s2}{.2f}\PY{l+s+si}{\PYZcb{}}\PY{l+s+s2}{\PYZdq{}}\PY{p}{)}\PY{p}{)}

    \PY{k}{return} \PY{n}{x\PYZus{}g}\PY{p}{,} \PY{n}{s2\PYZus{}g}
\end{Verbatim}
\end{tcolorbox}

    \begin{tcolorbox}[breakable, size=fbox, boxrule=1pt, pad at break*=1mm,colback=cellbackground, colframe=cellborder]
\prompt{In}{incolor}{16}{\boxspacing}
\begin{Verbatim}[commandchars=\\\{\}]
\PY{n}{moments}\PY{p}{(}\PY{n}{X1}\PY{p}{,} \PY{l+s+s2}{\PYZdq{}}\PY{l+s+s2}{X1}\PY{l+s+s2}{\PYZdq{}}\PY{p}{)}
\end{Verbatim}
\end{tcolorbox}

    \begin{Verbatim}[commandchars=\\\{\}]
Интервалы для X1:
1: [28.81, 34.05] | частота = 6 | середина = 31.43 | n\_k * X\_k = 188.58
2: [34.05, 39.29] | частота = 15 | середина = 36.67 | n\_k * X\_k = 550.03
3: [39.29, 44.53] | частота = 31 | середина = 41.91 | n\_k * X\_k = 1299.12
4: [44.53, 49.77] | частота = 42 | середина = 47.15 | n\_k * X\_k = 1980.14
5: [49.77, 55.00] | частота = 47 | середина = 52.38 | n\_k * X\_k = 2462.09
6: [55.00, 60.24] | частота = 33 | середина = 57.62 | n\_k * X\_k = 1901.59
7: [60.24, 65.48] | частота = 13 | середина = 62.86 | n\_k * X\_k = 817.22
8: [65.48, 70.72] | частота = 7 | середина = 68.10 | n\_k * X\_k = 476.71
9: [70.72, 75.96] | частота = 6 | середина = 73.34 | n\_k * X\_k = 440.04
Итог сумм: частот = 200, n\_k * X\_k = 10115.52
    \end{Verbatim}

    $\displaystyle \hat{X}_g = \frac{10115.52}{200} = 50.58$

    
    \begin{Verbatim}[commandchars=\\\{\}]

1: [28.81, 34.05] | частота = 6 | середина = 31.43 | n\_k * (X\_k - X\_g)\^{}2 =
2199.91
2: [34.05, 39.29] | частота = 15 | середина = 36.67 | n\_k * (X\_k - X\_g)\^{}2 =
2902.01
3: [39.29, 44.53] | частота = 31 | середина = 41.91 | n\_k * (X\_k - X\_g)\^{}2 =
2330.43
4: [44.53, 49.77] | частота = 42 | середина = 47.15 | n\_k * (X\_k - X\_g)\^{}2 =
494.55
5: [49.77, 55.00] | частота = 47 | середина = 52.38 | n\_k * (X\_k - X\_g)\^{}2 =
153.54
6: [55.00, 60.24] | частота = 33 | середина = 57.62 | n\_k * (X\_k - X\_g)\^{}2 =
1638.46
7: [60.24, 65.48] | частота = 13 | середина = 62.86 | n\_k * (X\_k - X\_g)\^{}2 =
1962.04
8: [65.48, 70.72] | частота = 7 | середина = 68.10 | n\_k * (X\_k - X\_g)\^{}2 =
2149.65
9: [70.72, 75.96] | частота = 6 | середина = 73.34 | n\_k * (X\_k - X\_g)\^{}2 =
3108.92
    \end{Verbatim}

    $\displaystyle \hat{\sigma}_g^2 = \frac{16939.51}{199} = 85.12$

    
            \begin{tcolorbox}[breakable, size=fbox, boxrule=.5pt, pad at break*=1mm, opacityfill=0]
\prompt{Out}{outcolor}{16}{\boxspacing}
\begin{Verbatim}[commandchars=\\\{\}]
(np.float64(50.57758333333332), np.float64(85.12315226673488))
\end{Verbatim}
\end{tcolorbox}
        
    \begin{tcolorbox}[breakable, size=fbox, boxrule=1pt, pad at break*=1mm,colback=cellbackground, colframe=cellborder]
\prompt{In}{incolor}{17}{\boxspacing}
\begin{Verbatim}[commandchars=\\\{\}]
\PY{n}{moments}\PY{p}{(}\PY{n}{X2}\PY{p}{,} \PY{l+s+s2}{\PYZdq{}}\PY{l+s+s2}{X2}\PY{l+s+s2}{\PYZdq{}}\PY{p}{)}
\end{Verbatim}
\end{tcolorbox}

    \begin{Verbatim}[commandchars=\\\{\}]
Интервалы для X2:
1: [26.80, 31.89] | частота = 18 | середина = 29.34 | n\_k * X\_k = 528.17
2: [31.89, 36.97] | частота = 26 | середина = 34.43 | n\_k * X\_k = 895.14
3: [36.97, 42.06] | частота = 28 | середина = 39.51 | n\_k * X\_k = 1106.39
4: [42.06, 47.14] | частота = 22 | середина = 44.60 | n\_k * X\_k = 981.19
5: [47.14, 52.23] | частота = 21 | середина = 49.68 | n\_k * X\_k = 1043.38
6: [52.23, 57.31] | частота = 26 | середина = 54.77 | n\_k * X\_k = 1424.03
7: [57.31, 62.40] | частота = 19 | середина = 59.86 | n\_k * X\_k = 1137.27
8: [62.40, 67.48] | частота = 18 | середина = 64.94 | n\_k * X\_k = 1168.95
9: [67.48, 72.57] | частота = 22 | середина = 70.03 | n\_k * X\_k = 1540.60
Итог сумм: частот = 200, n\_k * X\_k = 9825.12
    \end{Verbatim}

    $\displaystyle \hat{X}_g = \frac{9825.12}{200} = 49.13$

    
    \begin{Verbatim}[commandchars=\\\{\}]

1: [26.80, 31.89] | частота = 18 | середина = 29.34 | n\_k * (X\_k - X\_g)\^{}2 =
7044.47
2: [31.89, 36.97] | частота = 26 | середина = 34.43 | n\_k * (X\_k - X\_g)\^{}2 =
5616.24
3: [36.97, 42.06] | частота = 28 | середина = 39.51 | n\_k * (X\_k - X\_g)\^{}2 =
2586.77
4: [42.06, 47.14] | частота = 22 | середина = 44.60 | n\_k * (X\_k - X\_g)\^{}2 =
450.69
5: [47.14, 52.23] | частота = 21 | середина = 49.68 | n\_k * (X\_k - X\_g)\^{}2 = 6.57
6: [52.23, 57.31] | частота = 26 | середина = 54.77 | n\_k * (X\_k - X\_g)\^{}2 =
828.51
7: [57.31, 62.40] | частота = 19 | середина = 59.86 | n\_k * (X\_k - X\_g)\^{}2 =
2187.74
8: [62.40, 67.48] | частота = 18 | середина = 64.94 | n\_k * (X\_k - X\_g)\^{}2 =
4502.67
9: [67.48, 72.57] | частота = 22 | середина = 70.03 | n\_k * (X\_k - X\_g)\^{}2 =
9611.32
    \end{Verbatim}

    $\displaystyle \hat{\sigma}_g^2 = \frac{32834.99}{199} = 165.00$

    
            \begin{tcolorbox}[breakable, size=fbox, boxrule=.5pt, pad at break*=1mm, opacityfill=0]
\prompt{Out}{outcolor}{17}{\boxspacing}
\begin{Verbatim}[commandchars=\\\{\}]
(np.float64(49.12558888888888), np.float64(164.99994590123447))
\end{Verbatim}
\end{tcolorbox}
        
    \begin{tcolorbox}[breakable, size=fbox, boxrule=1pt, pad at break*=1mm,colback=cellbackground, colframe=cellborder]
\prompt{In}{incolor}{18}{\boxspacing}
\begin{Verbatim}[commandchars=\\\{\}]
\PY{n}{moments}\PY{p}{(}\PY{n}{X3}\PY{p}{,} \PY{l+s+s2}{\PYZdq{}}\PY{l+s+s2}{X3}\PY{l+s+s2}{\PYZdq{}}\PY{p}{)}
\end{Verbatim}
\end{tcolorbox}

    \begin{Verbatim}[commandchars=\\\{\}]
Интервалы для X3:
1: [22.47, 65.19] | частота = 104 | середина = 43.83 | n\_k * X\_k = 4558.55
2: [65.19, 107.92] | частота = 51 | середина = 86.56 | n\_k * X\_k = 4414.39
3: [107.92, 150.64] | частота = 23 | середина = 129.28 | n\_k * X\_k = 2973.47
4: [150.64, 193.37] | частота = 10 | середина = 172.01 | n\_k * X\_k = 1720.06
5: [193.37, 236.09] | частота = 6 | середина = 214.73 | n\_k * X\_k = 1288.38
6: [236.09, 278.82] | частота = 5 | середина = 257.45 | n\_k * X\_k = 1287.27
7: [278.82, 321.54] | частота = 0 | середина = 300.18 | n\_k * X\_k = 0.00
8: [321.54, 364.27] | частота = 0 | середина = 342.90 | n\_k * X\_k = 0.00
9: [364.27, 406.99] | частота = 1 | середина = 385.63 | n\_k * X\_k = 385.63
Итог сумм: частот = 200, n\_k * X\_k = 16627.74
    \end{Verbatim}

    $\displaystyle \hat{X}_g = \frac{16627.74}{200} = 83.14$

    
    \begin{Verbatim}[commandchars=\\\{\}]

1: [22.47, 65.19] | частота = 104 | середина = 43.83 | n\_k * (X\_k - X\_g)\^{}2 =
160680.01
2: [65.19, 107.92] | частота = 51 | середина = 86.56 | n\_k * (X\_k - X\_g)\^{}2 =
595.80
3: [107.92, 150.64] | частота = 23 | середина = 129.28 | n\_k * (X\_k - X\_g)\^{}2 =
48969.78
4: [150.64, 193.37] | частота = 10 | середина = 172.01 | n\_k * (X\_k - X\_g)\^{}2 =
78973.16
5: [193.37, 236.09] | частота = 6 | середина = 214.73 | n\_k * (X\_k - X\_g)\^{}2 =
103897.60
6: [236.09, 278.82] | частота = 5 | середина = 257.45 | n\_k * (X\_k - X\_g)\^{}2 =
151929.87
7: [278.82, 321.54] | частота = 0 | середина = 300.18 | n\_k * (X\_k - X\_g)\^{}2 =
0.00
8: [321.54, 364.27] | частота = 0 | середина = 342.90 | n\_k * (X\_k - X\_g)\^{}2 =
0.00
9: [364.27, 406.99] | частота = 1 | середина = 385.63 | n\_k * (X\_k - X\_g)\^{}2 =
91499.64
    \end{Verbatim}

    $\displaystyle \hat{\sigma}_g^2 = \frac{636545.87}{199} = 3198.72$

    
            \begin{tcolorbox}[breakable, size=fbox, boxrule=.5pt, pad at break*=1mm, opacityfill=0]
\prompt{Out}{outcolor}{18}{\boxspacing}
\begin{Verbatim}[commandchars=\\\{\}]
(np.float64(83.13871111111109), np.float64(3198.7229625341515))
\end{Verbatim}
\end{tcolorbox}
        
    Сравниваем с исходными оценками

\textbf{Для X1:}

{\def\LTcaptype{none} % do not increment counter
\begin{longtable}[]{@{}
  >{\raggedright\arraybackslash}p{(\linewidth - 6\tabcolsep) * \real{0.3738}}
  >{\raggedright\arraybackslash}p{(\linewidth - 6\tabcolsep) * \real{0.1869}}
  >{\raggedright\arraybackslash}p{(\linewidth - 6\tabcolsep) * \real{0.2523}}
  >{\raggedright\arraybackslash}p{(\linewidth - 6\tabcolsep) * \real{0.1869}}@{}}
\toprule\noalign{}
\begin{minipage}[b]{\linewidth}\raggedright
Показатель
\end{minipage} & \begin{minipage}[b]{\linewidth}\raggedright
По исходным данным
\end{minipage} & \begin{minipage}[b]{\linewidth}\raggedright
По сгруппированным данным
\end{minipage} & \begin{minipage}[b]{\linewidth}\raggedright
Дельта расхождения
\end{minipage} \\
\midrule\noalign{}
\endhead
\bottomrule\noalign{}
\endlastfoot
Среднее \(\bar{x}\) & 50.57 & 50.58 & 0.01 \\
Несмещенная дисперсия \(\hat{\sigma}^2\) & 86.77 & 85.12 & 1.65 \\
\end{longtable}
}

\textbf{Для X2:}

{\def\LTcaptype{none} % do not increment counter
\begin{longtable}[]{@{}
  >{\raggedright\arraybackslash}p{(\linewidth - 6\tabcolsep) * \real{0.3704}}
  >{\raggedright\arraybackslash}p{(\linewidth - 6\tabcolsep) * \real{0.1944}}
  >{\raggedright\arraybackslash}p{(\linewidth - 6\tabcolsep) * \real{0.2500}}
  >{\raggedright\arraybackslash}p{(\linewidth - 6\tabcolsep) * \real{0.1852}}@{}}
\toprule\noalign{}
\begin{minipage}[b]{\linewidth}\raggedright
Показатель
\end{minipage} & \begin{minipage}[b]{\linewidth}\raggedright
По исходным данным
\end{minipage} & \begin{minipage}[b]{\linewidth}\raggedright
По сгруппированным данным
\end{minipage} & \begin{minipage}[b]{\linewidth}\raggedright
Дельта расхождения
\end{minipage} \\
\midrule\noalign{}
\endhead
\bottomrule\noalign{}
\endlastfoot
Среднее \(\bar{x}\) & 49.10 & 49.13 & 0.03 \\
Несмещенная дисперсия \(\hat{\sigma}^2\) & 167.58 & 165.00 & 2.58 \\
\end{longtable}
}

\textbf{Для X3:}

{\def\LTcaptype{none} % do not increment counter
\begin{longtable}[]{@{}
  >{\raggedright\arraybackslash}p{(\linewidth - 6\tabcolsep) * \real{0.3670}}
  >{\raggedright\arraybackslash}p{(\linewidth - 6\tabcolsep) * \real{0.2018}}
  >{\raggedright\arraybackslash}p{(\linewidth - 6\tabcolsep) * \real{0.2477}}
  >{\raggedright\arraybackslash}p{(\linewidth - 6\tabcolsep) * \real{0.1835}}@{}}
\toprule\noalign{}
\begin{minipage}[b]{\linewidth}\raggedright
Показатель
\end{minipage} & \begin{minipage}[b]{\linewidth}\raggedright
По исходным данным
\end{minipage} & \begin{minipage}[b]{\linewidth}\raggedright
По сгруппированным данным
\end{minipage} & \begin{minipage}[b]{\linewidth}\raggedright
Дельта расхождения
\end{minipage} \\
\midrule\noalign{}
\endhead
\bottomrule\noalign{}
\endlastfoot
Среднее \(\bar{x}\) & 79.81 & 83.14 & 3.33 \\
Несмещенная дисперсия \(\hat{\sigma}^2\) & 3318.17 & 3198.72 & 119.44 \\
\end{longtable}
}

    \paragraph{Выводы}\label{ux432ux44bux432ux43eux434ux44b}

Оценки математического ожидания и дисперсии, полученные по
сгруппированным данным, оказались близки к оценкам по исходным данным
для всех трех столбцов. Это подтверждает, что группировка данных с
использованием правила Стерджеса (\(k = 9\) интервалов) позволяет
сохранить основные статистические характеристики выборки с приемлемой
точностью

\textbf{Для X1}

Расхождение между оценками минимально: среднее отличается на 0.01,
дисперсия - на 1.65. Столь малое расхождение свидетельствует о том, что
группировка практически не приводит к потере информации.

\textbf{Для X2}

Расхождение также мало: среднее отличается на 0.03, дисперсия - на 2.58.
Для равномерной выборки это ожидаемо, так как данные сравнительно ровно
заполняют интервал значений.

\textbf{Для X3}

Наблюдается наибольшее расхождение среди всех столбцов: среднее
отличается на 3.33, дисперсия - на 119.44. Это объясняется длинным
правым хвостом и наличием крупных значений, которые при группировке
частично сглаживаются.

    \subsubsection{6. Доверительные
интервалы}\label{ux434ux43eux432ux435ux440ux438ux442ux435ux43bux44cux43dux44bux435-ux438ux43dux442ux435ux440ux432ux430ux43bux44b}

Уровень доверия \(1 - \alpha = 0.95 \Rightarrow \alpha = 0.05\)

\textbf{Асимптотический ДИ для \(EX\)}

\[
EX \in \left(\bar{x} - z_{1 - \frac{\alpha}{2}}\frac{\hat{\sigma}}{\sqrt{n}}, \ \bar{x} + z_{1 - \frac{\alpha}{2}}\frac{\hat{\sigma}}{\sqrt{n}}\right)
\]

\(z_{1 - \frac{\alpha}{2}} = 1.96\) - квантиль стандартного нормального
распределения

\(\hat{\sigma}\) - несмещенное стандартное отклонение

    \begin{tcolorbox}[breakable, size=fbox, boxrule=1pt, pad at break*=1mm,colback=cellbackground, colframe=cellborder]
\prompt{In}{incolor}{19}{\boxspacing}
\begin{Verbatim}[commandchars=\\\{\}]
\PY{n+nb}{print}\PY{p}{(}\PY{n+nb}{round}\PY{p}{(}\PY{n}{stats}\PY{o}{.}\PY{n}{norm}\PY{o}{.}\PY{n}{ppf}\PY{p}{(}\PY{l+m+mi}{1} \PY{o}{\PYZhy{}} \PY{l+m+mf}{0.05} \PY{o}{/} \PY{l+m+mi}{2}\PY{p}{)}\PY{p}{,} \PY{l+m+mi}{2}\PY{p}{)}\PY{p}{)}
\end{Verbatim}
\end{tcolorbox}

    \begin{Verbatim}[commandchars=\\\{\}]
1.96
    \end{Verbatim}

    \textbf{Для X1 - нормальное распределение}

\begin{itemize}
\tightlist
\item
  \(\bar{x} = 50.57\)
\item
  \(\hat{\sigma} = 9.315\)
\item
  \(\frac{\hat{\sigma}}{\sqrt{n}} = \frac{9.315}{\sqrt{200}} = 0.659\)
\item
  Нижняя граница = \(50.57 - 1.96 \cdot 0.659 = 49.278\)
\item
  Верхняя граница = \(50.57 + 1.96 \cdot 0.659 = 51.860\)
\end{itemize}

\textbf{Таким образом, асимптотический 95\% ДИ для EX (X1):
\((49.278; 51.860)\)}

\begin{center}\rule{0.5\linewidth}{0.5pt}\end{center}

\textbf{Для X2 - равномерное распределение}

\begin{itemize}
\tightlist
\item
  \(\bar{x} = 49.10\)
\item
  \(\hat{\sigma} = 12.945\)
\item
  \(\frac{\hat{\sigma}}{\sqrt{n}} = \frac{12.945}{\sqrt{200}} = 0.915\)
\item
  Нижняя граница = \(49.10 - 1.96 \cdot 0.915 = 47.305\)
\item
  Верхняя граница = \(49.10 + 1.96 \cdot 0.915 = 50.894\)
\end{itemize}

\textbf{Таким образом, асимптотический 95\% ДИ для EX (X2):
\((47.305; 50.894)\)}

\begin{center}\rule{0.5\linewidth}{0.5pt}\end{center}

\textbf{Для X3 - экспоненциальное распределение}

\begin{itemize}
\tightlist
\item
  \(\bar{x} = 79.81\)
\item
  \(\hat{\sigma} = 57.604\)
\item
  \(\frac{\hat{\sigma}}{\sqrt{n}} = \frac{57.604}{\sqrt{200}} = 4.073\)
\item
  Нижняя граница = \(79.81 - 1.96 \cdot 4.073 = 71.826\)
\item
  Верхняя граница = \(79.81 + 1.96 \cdot 4.073 = 87.792\)
\end{itemize}

\textbf{Таким образом, асимптотический 95\% ДИ для EX (X3):
\((71.826; 87.792)\)}

    \textbf{Точные ДИ для параметров нормального распределения}

\(N(\mu, \sigma^2) \Rightarrow\)

\[
\mu \in \left(\bar{x} - t_{1 - \frac{\alpha}{2}, n - 1}\frac{\hat{\sigma}}{\sqrt{n}}, \ \bar{x} + t_{1 - \frac{\alpha}{2}, n - 1}\frac{\hat{\sigma}}{\sqrt{n}}\right),
\quad
\sigma^2 \in \left(\frac{(n-1)\hat{\sigma}^2}{\chi^2_{1 - \frac{\alpha}{2}, n - 1}}, \ \frac{(n-1)\hat{\sigma}^2}{\chi^2_{\frac{\alpha}{2}, n - 1}}\right)
\]

Где \(t_{1 - \frac{\alpha}{2}, n - 1}\) - квантиль уровня
\(1 - \frac{\alpha}{2}\) распределения Стьюдента с \(n - 1\) степенью
свободы

\(\chi^2_{1 - \frac{\alpha}{2}, n - 1}\) и
\(\chi^2_{\frac{\alpha}{2}, n - 1}\) - квантили распределения \(\chi^2\)
с \(n - 1\) степенью свободы.

    \begin{tcolorbox}[breakable, size=fbox, boxrule=1pt, pad at break*=1mm,colback=cellbackground, colframe=cellborder]
\prompt{In}{incolor}{20}{\boxspacing}
\begin{Verbatim}[commandchars=\\\{\}]
\PY{c+c1}{\PYZsh{} Считаем нужный квантиль распределения Стьюдента}

\PY{n}{t\PYZus{}value} \PY{o}{=} \PY{n}{stats}\PY{o}{.}\PY{n}{t}\PY{o}{.}\PY{n}{ppf}\PY{p}{(}\PY{l+m+mi}{1} \PY{o}{\PYZhy{}} \PY{l+m+mf}{0.05} \PY{o}{/} \PY{l+m+mi}{2}\PY{p}{,} \PY{n}{df}\PY{o}{=}\PY{n}{n} \PY{o}{\PYZhy{}} \PY{l+m+mi}{1}\PY{p}{)}
\PY{n+nb}{print}\PY{p}{(}\PY{n}{t\PYZus{}value}\PY{p}{)}
\end{Verbatim}
\end{tcolorbox}

    \begin{Verbatim}[commandchars=\\\{\}]
1.9719565442517535
    \end{Verbatim}

    \begin{tcolorbox}[breakable, size=fbox, boxrule=1pt, pad at break*=1mm,colback=cellbackground, colframe=cellborder]
\prompt{In}{incolor}{21}{\boxspacing}
\begin{Verbatim}[commandchars=\\\{\}]
\PY{c+c1}{\PYZsh{} Считаем нижний и верхний квантили для распределения хи\PYZhy{}квадрат}

\PY{n}{chi2\PYZus{}lower} \PY{o}{=} \PY{n}{stats}\PY{o}{.}\PY{n}{chi2}\PY{o}{.}\PY{n}{ppf}\PY{p}{(}\PY{l+m+mf}{0.05} \PY{o}{/} \PY{l+m+mi}{2}\PY{p}{,} \PY{n}{df}\PY{o}{=}\PY{n}{n} \PY{o}{\PYZhy{}} \PY{l+m+mi}{1}\PY{p}{)}

\PY{n}{chi2\PYZus{}upper} \PY{o}{=} \PY{n}{stats}\PY{o}{.}\PY{n}{chi2}\PY{o}{.}\PY{n}{ppf}\PY{p}{(}\PY{l+m+mi}{1} \PY{o}{\PYZhy{}} \PY{l+m+mf}{0.05} \PY{o}{/} \PY{l+m+mi}{2}\PY{p}{,} \PY{n}{df}\PY{o}{=}\PY{n}{n} \PY{o}{\PYZhy{}} \PY{l+m+mi}{1}\PY{p}{)}

\PY{n}{display}\PY{p}{(}\PY{n}{Math}\PY{p}{(}
    \PY{l+s+sa}{rf}\PY{l+s+s2}{\PYZdq{}}\PY{l+s+s2}{\PYZbs{}}\PY{l+s+s2}{chi\PYZca{}2\PYZus{}}\PY{l+s+se}{\PYZob{}\PYZob{}}\PY{l+s+s2}{1 \PYZhy{} }\PY{l+s+s2}{\PYZbs{}}\PY{l+s+s2}{frac}\PY{l+s+se}{\PYZob{}\PYZob{}}\PY{l+s+s2}{\PYZbs{}}\PY{l+s+s2}{alpha}\PY{l+s+se}{\PYZcb{}\PYZcb{}}\PY{l+s+se}{\PYZob{}\PYZob{}}\PY{l+s+s2}{2}\PY{l+s+se}{\PYZcb{}\PYZcb{}}\PY{l+s+s2}{, n \PYZhy{} 1}\PY{l+s+se}{\PYZcb{}\PYZcb{}}\PY{l+s+s2}{ = }\PY{l+s+si}{\PYZob{}}\PY{n}{chi2\PYZus{}upper}\PY{l+s+si}{\PYZcb{}}\PY{l+s+s2}{, }\PY{l+s+s2}{\PYZbs{}}\PY{l+s+s2}{ }\PY{l+s+s2}{\PYZbs{}}\PY{l+s+s2}{chi\PYZca{}2\PYZus{}}\PY{l+s+se}{\PYZob{}\PYZob{}}\PY{l+s+s2}{\PYZbs{}}\PY{l+s+s2}{frac}\PY{l+s+se}{\PYZob{}\PYZob{}}\PY{l+s+s2}{\PYZbs{}}\PY{l+s+s2}{alpha}\PY{l+s+se}{\PYZcb{}\PYZcb{}}\PY{l+s+se}{\PYZob{}\PYZob{}}\PY{l+s+s2}{2}\PY{l+s+se}{\PYZcb{}\PYZcb{}}\PY{l+s+s2}{, n \PYZhy{} 1}\PY{l+s+se}{\PYZcb{}\PYZcb{}}\PY{l+s+s2}{ = }\PY{l+s+si}{\PYZob{}}\PY{n}{chi2\PYZus{}lower}\PY{l+s+si}{\PYZcb{}}\PY{l+s+s2}{\PYZdq{}}\PY{p}{)}\PY{p}{)}
\end{Verbatim}
\end{tcolorbox}

    $\displaystyle \chi^2_{1 - \frac{\alpha}{2}, n - 1} = 239.9596818276442, \ \chi^2_{\frac{\alpha}{2}, n - 1} = 161.82618239364686$

    
    Получив все данные, посчитаем границы для \(\mu\)

\begin{itemize}
\tightlist
\item
  \(\bar{x} = 50.57\)
\item
  \(\hat{\sigma} = 9.315\)
\item
  \(\frac{\hat{\sigma}}{\sqrt{n}} = \frac{9.315}{\sqrt{200}} = 0.659\)
\end{itemize}

Нижняя граница = \(50.57 - 1.972 \cdot 0.659 = 49.270\)

Верхняя граница = \(50.57 + 1.972 \cdot 0.659 = 51.868\)

\textbf{Таким образом, получаем, что точный 95\% ДИ для \(\mu\)
нормального распределения X1:} \((49.270, \ 51.868)\)

\begin{center}\rule{0.5\linewidth}{0.5pt}\end{center}

Теперь считаем границы для \(\sigma^2\)

\begin{itemize}
\tightlist
\item
  \(\hat{\sigma}^2 = 86.77\)
\item
  \(n - 1 = 199\)
\end{itemize}

Нижняя граница = \(\frac{199 \cdot 86.77}{239.96} = 71.959\)

Верхняя граница = \(\frac{199 \cdot 86.77}{161.83} = 106.703\)

\textbf{Таким образом, получаем, что точный 95\% ДИ для \(\sigma^2\)
нормального распределения X1:} \((71.959, \ 106.703)\)

    Интерпретация

Доверительный интервал - это интервальная оценка, которая с заданным
уровнем доверия накрывает истинное значение параметра генеральной
совокупности. Важно понимать, что интервал является случайным, так как
он построен по случайной выборке, а истинный параметр --- фиксированное,
но неизвестное число. Поэтому корректная формулировка: ``Если
многократно извлекать выборки и строить доверительные интервалы, то
примерно в 95\% случаев они будут содержать истинный параметр''.
Некорректно утверждать, что ``вероятность того, что параметр лежит в
интервале, равна 0.95'', поскольку параметр не является случайной
величиной

    \subsubsection{7. Итоги и
выводы}\label{ux438ux442ux43eux433ux438-ux438-ux432ux44bux432ux43eux434ux44b}

В результате выполнения работы для трех выборок объемом \(n = 200\) были
определены вероятностные законы распределения, оценены параметры и
построены доверительные интервалы.

Для X1 выбрана модель нормального распределения \(N(a, \sigma)\) с
оценками параметров \(a = 50.57\), \(\sigma = 9.29\) (метод моментов и
ММП дают одинаковый результат). 95\% доверительный интервал для
математического ожидания составил \((49.278; 51.860)\), что указывает на
хорошую точность оценки. Данные близки к симметричным, гистограмма имеет
колоколообразную форму.

Для X2 выбрана модель равномерного распределения \(U(a, b)\) с оценками
\(a = 26.80\), \(b = 72.57\) по ММП. Доверительный интервал для \(EX\)
составил \((47.305; 50.894)\). По гистограмме видно, что значения
сравнительно равномерно заполняют конечный интервал без выраженных
хвостов и выбросов.

Для X3 выбрана модель экспоненциального распределения со сдвигом
\(Exp(\lambda, c)\) с оценками \(\lambda = 0.0174\), \(c = 22.47\)
(ММП). Доверительный интервал для \(EX\) составил \((71.826; 87.792)\),
ширина интервала заметно больше, чем у X1 и X2, что отражает большой
разброс данных и тяжелый правый хвост.

Практический вывод: для X1 уместна нормальная модель, для X2 -
равномерная, для X3 - экспоненциальная со сдвигом. Выбранные модели
согласуются с графиками, выборочными характеристиками и оценками
вероятностей.

    \subsubsection{8. Бонус}\label{ux431ux43eux43dux443ux441}

    \begin{tcolorbox}[breakable, size=fbox, boxrule=1pt, pad at break*=1mm,colback=cellbackground, colframe=cellborder]
\prompt{In}{incolor}{22}{\boxspacing}
\begin{Verbatim}[commandchars=\\\{\}]
\PY{c+c1}{\PYZsh{} Гистограмма с правилом Стерджеса}
\PY{n}{bins} \PY{o}{=} \PY{n+nb}{int}\PY{p}{(}\PY{l+m+mi}{1} \PY{o}{+} \PY{n}{np}\PY{o}{.}\PY{n}{log2}\PY{p}{(}\PY{n}{n}\PY{p}{)}\PY{p}{)}
\PY{n}{plt}\PY{o}{.}\PY{n}{figure}\PY{p}{(}\PY{n}{figsize}\PY{o}{=}\PY{p}{(}\PY{l+m+mi}{12}\PY{p}{,} \PY{l+m+mi}{5}\PY{p}{)}\PY{p}{)}

\PY{n}{plt}\PY{o}{.}\PY{n}{subplot}\PY{p}{(}\PY{l+m+mi}{1}\PY{p}{,} \PY{l+m+mi}{2}\PY{p}{,} \PY{l+m+mi}{1}\PY{p}{)}
\PY{n}{plt}\PY{o}{.}\PY{n}{hist}\PY{p}{(}\PY{n}{X4}\PY{p}{,} \PY{n}{bins}\PY{o}{=}\PY{n}{bins}\PY{p}{,} \PY{n}{density}\PY{o}{=}\PY{k+kc}{True}\PY{p}{,} \PY{n}{alpha}\PY{o}{=}\PY{l+m+mf}{0.7}\PY{p}{,} \PY{n}{edgecolor}\PY{o}{=}\PY{l+s+s1}{\PYZsq{}}\PY{l+s+s1}{black}\PY{l+s+s1}{\PYZsq{}}\PY{p}{,} \PY{n}{color}\PY{o}{=}\PY{l+s+s1}{\PYZsq{}}\PY{l+s+s1}{lightblue}\PY{l+s+s1}{\PYZsq{}}\PY{p}{)}
\PY{n}{plt}\PY{o}{.}\PY{n}{title}\PY{p}{(}\PY{l+s+s1}{\PYZsq{}}\PY{l+s+s1}{Гистограмма X4}\PY{l+s+s1}{\PYZsq{}}\PY{p}{)}
\PY{n}{plt}\PY{o}{.}\PY{n}{xlabel}\PY{p}{(}\PY{l+s+s1}{\PYZsq{}}\PY{l+s+s1}{Значение}\PY{l+s+s1}{\PYZsq{}}\PY{p}{)}
\PY{n}{plt}\PY{o}{.}\PY{n}{ylabel}\PY{p}{(}\PY{l+s+s1}{\PYZsq{}}\PY{l+s+s1}{Плотность}\PY{l+s+s1}{\PYZsq{}}\PY{p}{)}
\PY{n}{plt}\PY{o}{.}\PY{n}{grid}\PY{p}{(}\PY{k+kc}{True}\PY{p}{,} \PY{n}{alpha}\PY{o}{=}\PY{l+m+mf}{0.3}\PY{p}{)}

\PY{c+c1}{\PYZsh{} Эмпирическая функция распределения}
\PY{n}{ecdf} \PY{o}{=} \PY{n}{ECDF}\PY{p}{(}\PY{n}{X4}\PY{p}{)}

\PY{n}{plt}\PY{o}{.}\PY{n}{subplot}\PY{p}{(}\PY{l+m+mi}{1}\PY{p}{,} \PY{l+m+mi}{2}\PY{p}{,} \PY{l+m+mi}{2}\PY{p}{)}
\PY{n}{plt}\PY{o}{.}\PY{n}{step}\PY{p}{(}\PY{n}{ecdf}\PY{o}{.}\PY{n}{x}\PY{p}{,} \PY{n}{ecdf}\PY{o}{.}\PY{n}{y}\PY{p}{,} \PY{n}{where}\PY{o}{=}\PY{l+s+s1}{\PYZsq{}}\PY{l+s+s1}{post}\PY{l+s+s1}{\PYZsq{}}\PY{p}{)}
\PY{n}{plt}\PY{o}{.}\PY{n}{title}\PY{p}{(}\PY{l+s+s1}{\PYZsq{}}\PY{l+s+s1}{Эмпирическая функция распределения X4}\PY{l+s+s1}{\PYZsq{}}\PY{p}{)}
\PY{n}{plt}\PY{o}{.}\PY{n}{xlabel}\PY{p}{(}\PY{l+s+s1}{\PYZsq{}}\PY{l+s+s1}{x}\PY{l+s+s1}{\PYZsq{}}\PY{p}{)}
\PY{n}{plt}\PY{o}{.}\PY{n}{ylabel}\PY{p}{(}\PY{l+s+s1}{\PYZsq{}}\PY{l+s+s1}{Fn(x)}\PY{l+s+s1}{\PYZsq{}}\PY{p}{)}
\PY{n}{plt}\PY{o}{.}\PY{n}{grid}\PY{p}{(}\PY{k+kc}{True}\PY{p}{,} \PY{n}{alpha}\PY{o}{=}\PY{l+m+mf}{0.3}\PY{p}{)}

\PY{n}{plt}\PY{o}{.}\PY{n}{tight\PYZus{}layout}\PY{p}{(}\PY{p}{)}
\PY{n}{plt}\PY{o}{.}\PY{n}{show}\PY{p}{(}\PY{p}{)}

\PY{c+c1}{\PYZsh{} Числовые характеристики}
\PY{n+nb}{print}\PY{p}{(}\PY{l+s+s2}{\PYZdq{}}\PY{l+s+s2}{Числовые характеристики X4:}\PY{l+s+s2}{\PYZdq{}}\PY{p}{)}
\PY{n+nb}{print}\PY{p}{(}\PY{l+s+sa}{f}\PY{l+s+s2}{\PYZdq{}}\PY{l+s+s2}{Среднее: }\PY{l+s+si}{\PYZob{}}\PY{n}{np}\PY{o}{.}\PY{n}{mean}\PY{p}{(}\PY{n}{X4}\PY{p}{)}\PY{l+s+si}{:}\PY{l+s+s2}{.2f}\PY{l+s+si}{\PYZcb{}}\PY{l+s+s2}{\PYZdq{}}\PY{p}{)}
\PY{n+nb}{print}\PY{p}{(}\PY{l+s+sa}{f}\PY{l+s+s2}{\PYZdq{}}\PY{l+s+s2}{Медиана: }\PY{l+s+si}{\PYZob{}}\PY{n}{np}\PY{o}{.}\PY{n}{median}\PY{p}{(}\PY{n}{X4}\PY{p}{)}\PY{l+s+si}{:}\PY{l+s+s2}{.2f}\PY{l+s+si}{\PYZcb{}}\PY{l+s+s2}{\PYZdq{}}\PY{p}{)}
\PY{n+nb}{print}\PY{p}{(}\PY{l+s+sa}{f}\PY{l+s+s2}{\PYZdq{}}\PY{l+s+s2}{Стандартное отклонение: }\PY{l+s+si}{\PYZob{}}\PY{n}{np}\PY{o}{.}\PY{n}{std}\PY{p}{(}\PY{n}{X4}\PY{p}{,}\PY{+w}{ }\PY{n}{ddof}\PY{o}{=}\PY{l+m+mi}{1}\PY{p}{)}\PY{l+s+si}{:}\PY{l+s+s2}{.2f}\PY{l+s+si}{\PYZcb{}}\PY{l+s+s2}{\PYZdq{}}\PY{p}{)}
\PY{n+nb}{print}\PY{p}{(}\PY{l+s+sa}{f}\PY{l+s+s2}{\PYZdq{}}\PY{l+s+s2}{Минимум: }\PY{l+s+si}{\PYZob{}}\PY{n}{np}\PY{o}{.}\PY{n}{min}\PY{p}{(}\PY{n}{X4}\PY{p}{)}\PY{l+s+si}{:}\PY{l+s+s2}{.2f}\PY{l+s+si}{\PYZcb{}}\PY{l+s+s2}{\PYZdq{}}\PY{p}{)}
\PY{n+nb}{print}\PY{p}{(}\PY{l+s+sa}{f}\PY{l+s+s2}{\PYZdq{}}\PY{l+s+s2}{Максимум: }\PY{l+s+si}{\PYZob{}}\PY{n}{np}\PY{o}{.}\PY{n}{max}\PY{p}{(}\PY{n}{X4}\PY{p}{)}\PY{l+s+si}{:}\PY{l+s+s2}{.2f}\PY{l+s+si}{\PYZcb{}}\PY{l+s+s2}{\PYZdq{}}\PY{p}{)}
\end{Verbatim}
\end{tcolorbox}

    \begin{center}
    \adjustimage{max size={0.9\linewidth}{0.9\paperheight}}{RGR1_A-2_files/RGR1_A-2_54_0.png}
    \end{center}
    { \hspace*{\fill} \\}
    
    \begin{Verbatim}[commandchars=\\\{\}]
Числовые характеристики X4:
Среднее: 71.30
Медиана: 46.76
Стандартное отклонение: 37.40
Минимум: 22.80
Максимум: 126.84
    \end{Verbatim}

    Гистограмма X4 показывает четкую бимодальность - два явно выраженных
пика:

Первый пик: значения примерно 34.36-45.92 (плотность 0.032) Второй пик:
значения примерно 103.72-115.28 (плотность 0.022)

Между ними наблюдается провал: на интервале примерно 57-92 плотность
равна 0, что очень похоже на смесь двух разных групп наблюдений.

\textbf{Признаки неоднородности данных}

\begin{enumerate}
\def\labelenumi{\arabic{enumi}.}
\item
  Бимодальность гистограммы. Наличие двух разделенных пиков является
  главным признаком того, что выборка не является однородной и
  представляет собой смесь двух распределений.
\item
  Разрыв между группами. Отсутствие наблюдений в промежутке между двумя
  группами говорит о том, что данные естественным образом распадаются на
  два кластера без переходной зоны.
\item
  Расхождение среднего и медианы. Среднее (71.30) существенно отличается
  от медианы (46.76). Здесь это расхождение объясняется тем, что в
  выборке смешаны два режима.
\item
  Большое стандартное отклонение. Стандартное отклонение (37.40) велико
  по сравнению с разбросом внутри каждой из двух визуально выделяемых
  групп.
\item
  Форма эмпирической функции распределения. ЭФР имеет два участка
  быстрого роста: первый - в области низких значений, второй - в области
  высоких значений. Между ними наблюдается пологий участок, что
  характерно для смеси распределений.
\end{enumerate}

    \begin{tcolorbox}[breakable, size=fbox, boxrule=1pt, pad at break*=1mm,colback=cellbackground, colframe=cellborder]
\prompt{In}{incolor}{23}{\boxspacing}
\begin{Verbatim}[commandchars=\\\{\}]
\PY{c+c1}{\PYZsh{} Кластеризация}

\PY{c+c1}{\PYZsh{} Подготовка данных}
\PY{n}{X4\PYZus{}reshaped} \PY{o}{=} \PY{n}{X4}\PY{o}{.}\PY{n}{values}\PY{o}{.}\PY{n}{reshape}\PY{p}{(}\PY{o}{\PYZhy{}}\PY{l+m+mi}{1}\PY{p}{,} \PY{l+m+mi}{1}\PY{p}{)}

\PY{c+c1}{\PYZsh{} K\PYZhy{}means кластеризация}
\PY{n}{kmeans} \PY{o}{=} \PY{n}{KMeans}\PY{p}{(}\PY{n}{n\PYZus{}clusters}\PY{o}{=}\PY{l+m+mi}{2}\PY{p}{,} \PY{n}{random\PYZus{}state}\PY{o}{=}\PY{l+m+mi}{42}\PY{p}{,} \PY{n}{n\PYZus{}init}\PY{o}{=}\PY{l+m+mi}{10}\PY{p}{)}
\PY{n}{raw\PYZus{}labels} \PY{o}{=} \PY{n}{kmeans}\PY{o}{.}\PY{n}{fit\PYZus{}predict}\PY{p}{(}\PY{n}{X4\PYZus{}reshaped}\PY{p}{)}

\PY{c+c1}{\PYZsh{} Упорядочим метки так, чтобы Кластер 1 был с меньшими значениями}
\PY{n}{centers} \PY{o}{=} \PY{n}{kmeans}\PY{o}{.}\PY{n}{cluster\PYZus{}centers\PYZus{}}\PY{o}{.}\PY{n}{flatten}\PY{p}{(}\PY{p}{)}
\PY{n}{order} \PY{o}{=} \PY{n}{np}\PY{o}{.}\PY{n}{argsort}\PY{p}{(}\PY{n}{centers}\PY{p}{)}
\PY{n}{labels} \PY{o}{=} \PY{n}{np}\PY{o}{.}\PY{n}{where}\PY{p}{(}\PY{n}{raw\PYZus{}labels} \PY{o}{==} \PY{n}{order}\PY{p}{[}\PY{l+m+mi}{0}\PY{p}{]}\PY{p}{,} \PY{l+m+mi}{0}\PY{p}{,} \PY{l+m+mi}{1}\PY{p}{)}

\PY{c+c1}{\PYZsh{} Разделение на кластеры}
\PY{n}{cluster1} \PY{o}{=} \PY{n}{X4}\PY{p}{[}\PY{n}{labels} \PY{o}{==} \PY{l+m+mi}{0}\PY{p}{]}
\PY{n}{cluster2} \PY{o}{=} \PY{n}{X4}\PY{p}{[}\PY{n}{labels} \PY{o}{==} \PY{l+m+mi}{1}\PY{p}{]}

\PY{n+nb}{print}\PY{p}{(}\PY{l+s+s2}{\PYZdq{}}\PY{l+s+s2}{Кластеризация X4 (k\PYZhy{}means, k=2)}\PY{l+s+s2}{\PYZdq{}}\PY{p}{)}

\PY{n+nb}{print}\PY{p}{(}\PY{l+s+sa}{f}\PY{l+s+s2}{\PYZdq{}}\PY{l+s+se}{\PYZbs{}n}\PY{l+s+s2}{Кластер 1:}\PY{l+s+s2}{\PYZdq{}}\PY{p}{)}
\PY{n+nb}{print}\PY{p}{(}\PY{l+s+sa}{f}\PY{l+s+s2}{\PYZdq{}}\PY{l+s+s2}{Размер: }\PY{l+s+si}{\PYZob{}}\PY{n+nb}{len}\PY{p}{(}\PY{n}{cluster1}\PY{p}{)}\PY{l+s+si}{\PYZcb{}}\PY{l+s+s2}{ (}\PY{l+s+si}{\PYZob{}}\PY{n+nb}{len}\PY{p}{(}\PY{n}{cluster1}\PY{p}{)}\PY{+w}{ }\PY{o}{/}\PY{+w}{ }\PY{n+nb}{len}\PY{p}{(}\PY{n}{X4}\PY{p}{)}\PY{+w}{ }\PY{o}{*}\PY{+w}{ }\PY{l+m+mi}{100}\PY{l+s+si}{:}\PY{l+s+s2}{.1f}\PY{l+s+si}{\PYZcb{}}\PY{l+s+s2}{\PYZpc{})}\PY{l+s+s2}{\PYZdq{}}\PY{p}{)}
\PY{n+nb}{print}\PY{p}{(}\PY{l+s+sa}{f}\PY{l+s+s2}{\PYZdq{}}\PY{l+s+s2}{Среднее: }\PY{l+s+si}{\PYZob{}}\PY{n}{np}\PY{o}{.}\PY{n}{mean}\PY{p}{(}\PY{n}{cluster1}\PY{p}{)}\PY{l+s+si}{:}\PY{l+s+s2}{.2f}\PY{l+s+si}{\PYZcb{}}\PY{l+s+s2}{\PYZdq{}}\PY{p}{)}
\PY{n+nb}{print}\PY{p}{(}\PY{l+s+sa}{f}\PY{l+s+s2}{\PYZdq{}}\PY{l+s+s2}{Медиана: }\PY{l+s+si}{\PYZob{}}\PY{n}{np}\PY{o}{.}\PY{n}{median}\PY{p}{(}\PY{n}{cluster1}\PY{p}{)}\PY{l+s+si}{:}\PY{l+s+s2}{.2f}\PY{l+s+si}{\PYZcb{}}\PY{l+s+s2}{\PYZdq{}}\PY{p}{)}
\PY{n+nb}{print}\PY{p}{(}\PY{l+s+sa}{f}\PY{l+s+s2}{\PYZdq{}}\PY{l+s+s2}{Дисперсия: }\PY{l+s+si}{\PYZob{}}\PY{n}{np}\PY{o}{.}\PY{n}{var}\PY{p}{(}\PY{n}{cluster1}\PY{p}{,}\PY{+w}{ }\PY{n}{ddof}\PY{o}{=}\PY{l+m+mi}{1}\PY{p}{)}\PY{l+s+si}{:}\PY{l+s+s2}{.2f}\PY{l+s+si}{\PYZcb{}}\PY{l+s+s2}{\PYZdq{}}\PY{p}{)}
\PY{n+nb}{print}\PY{p}{(}\PY{l+s+sa}{f}\PY{l+s+s2}{\PYZdq{}}\PY{l+s+s2}{Минимум: }\PY{l+s+si}{\PYZob{}}\PY{n}{np}\PY{o}{.}\PY{n}{min}\PY{p}{(}\PY{n}{cluster1}\PY{p}{)}\PY{l+s+si}{:}\PY{l+s+s2}{.2f}\PY{l+s+si}{\PYZcb{}}\PY{l+s+s2}{\PYZdq{}}\PY{p}{)}
\PY{n+nb}{print}\PY{p}{(}\PY{l+s+sa}{f}\PY{l+s+s2}{\PYZdq{}}\PY{l+s+s2}{Максимум: }\PY{l+s+si}{\PYZob{}}\PY{n}{np}\PY{o}{.}\PY{n}{max}\PY{p}{(}\PY{n}{cluster1}\PY{p}{)}\PY{l+s+si}{:}\PY{l+s+s2}{.2f}\PY{l+s+si}{\PYZcb{}}\PY{l+s+s2}{\PYZdq{}}\PY{p}{)}

\PY{n+nb}{print}\PY{p}{(}\PY{l+s+sa}{f}\PY{l+s+s2}{\PYZdq{}}\PY{l+s+se}{\PYZbs{}n}\PY{l+s+s2}{Кластер 2:}\PY{l+s+s2}{\PYZdq{}}\PY{p}{)}
\PY{n+nb}{print}\PY{p}{(}\PY{l+s+sa}{f}\PY{l+s+s2}{\PYZdq{}}\PY{l+s+s2}{Размер: }\PY{l+s+si}{\PYZob{}}\PY{n+nb}{len}\PY{p}{(}\PY{n}{cluster2}\PY{p}{)}\PY{l+s+si}{\PYZcb{}}\PY{l+s+s2}{ (}\PY{l+s+si}{\PYZob{}}\PY{n+nb}{len}\PY{p}{(}\PY{n}{cluster2}\PY{p}{)}\PY{+w}{ }\PY{o}{/}\PY{+w}{ }\PY{n+nb}{len}\PY{p}{(}\PY{n}{X4}\PY{p}{)}\PY{+w}{ }\PY{o}{*}\PY{+w}{ }\PY{l+m+mi}{100}\PY{l+s+si}{:}\PY{l+s+s2}{.1f}\PY{l+s+si}{\PYZcb{}}\PY{l+s+s2}{\PYZpc{})}\PY{l+s+s2}{\PYZdq{}}\PY{p}{)}
\PY{n+nb}{print}\PY{p}{(}\PY{l+s+sa}{f}\PY{l+s+s2}{\PYZdq{}}\PY{l+s+s2}{Среднее: }\PY{l+s+si}{\PYZob{}}\PY{n}{np}\PY{o}{.}\PY{n}{mean}\PY{p}{(}\PY{n}{cluster2}\PY{p}{)}\PY{l+s+si}{:}\PY{l+s+s2}{.2f}\PY{l+s+si}{\PYZcb{}}\PY{l+s+s2}{\PYZdq{}}\PY{p}{)}
\PY{n+nb}{print}\PY{p}{(}\PY{l+s+sa}{f}\PY{l+s+s2}{\PYZdq{}}\PY{l+s+s2}{Медиана: }\PY{l+s+si}{\PYZob{}}\PY{n}{np}\PY{o}{.}\PY{n}{median}\PY{p}{(}\PY{n}{cluster2}\PY{p}{)}\PY{l+s+si}{:}\PY{l+s+s2}{.2f}\PY{l+s+si}{\PYZcb{}}\PY{l+s+s2}{\PYZdq{}}\PY{p}{)}
\PY{n+nb}{print}\PY{p}{(}\PY{l+s+sa}{f}\PY{l+s+s2}{\PYZdq{}}\PY{l+s+s2}{Дисперсия: }\PY{l+s+si}{\PYZob{}}\PY{n}{np}\PY{o}{.}\PY{n}{var}\PY{p}{(}\PY{n}{cluster2}\PY{p}{,}\PY{+w}{ }\PY{n}{ddof}\PY{o}{=}\PY{l+m+mi}{1}\PY{p}{)}\PY{l+s+si}{:}\PY{l+s+s2}{.2f}\PY{l+s+si}{\PYZcb{}}\PY{l+s+s2}{\PYZdq{}}\PY{p}{)}
\PY{n+nb}{print}\PY{p}{(}\PY{l+s+sa}{f}\PY{l+s+s2}{\PYZdq{}}\PY{l+s+s2}{Минимум: }\PY{l+s+si}{\PYZob{}}\PY{n}{np}\PY{o}{.}\PY{n}{min}\PY{p}{(}\PY{n}{cluster2}\PY{p}{)}\PY{l+s+si}{:}\PY{l+s+s2}{.2f}\PY{l+s+si}{\PYZcb{}}\PY{l+s+s2}{\PYZdq{}}\PY{p}{)}
\PY{n+nb}{print}\PY{p}{(}\PY{l+s+sa}{f}\PY{l+s+s2}{\PYZdq{}}\PY{l+s+s2}{Максимум: }\PY{l+s+si}{\PYZob{}}\PY{n}{np}\PY{o}{.}\PY{n}{max}\PY{p}{(}\PY{n}{cluster2}\PY{p}{)}\PY{l+s+si}{:}\PY{l+s+s2}{.2f}\PY{l+s+si}{\PYZcb{}}\PY{l+s+s2}{\PYZdq{}}\PY{p}{)}

\PY{n+nb}{print}\PY{p}{(}\PY{l+s+sa}{f}\PY{l+s+s2}{\PYZdq{}}\PY{l+s+se}{\PYZbs{}n}\PY{l+s+s2}{Общая выборка X4:}\PY{l+s+s2}{\PYZdq{}}\PY{p}{)}
\PY{n+nb}{print}\PY{p}{(}\PY{l+s+sa}{f}\PY{l+s+s2}{\PYZdq{}}\PY{l+s+s2}{Среднее: }\PY{l+s+si}{\PYZob{}}\PY{n}{np}\PY{o}{.}\PY{n}{mean}\PY{p}{(}\PY{n}{X4}\PY{p}{)}\PY{l+s+si}{:}\PY{l+s+s2}{.2f}\PY{l+s+si}{\PYZcb{}}\PY{l+s+s2}{\PYZdq{}}\PY{p}{)}
\PY{n+nb}{print}\PY{p}{(}\PY{l+s+sa}{f}\PY{l+s+s2}{\PYZdq{}}\PY{l+s+s2}{Дисперсия: }\PY{l+s+si}{\PYZob{}}\PY{n}{np}\PY{o}{.}\PY{n}{var}\PY{p}{(}\PY{n}{X4}\PY{p}{,}\PY{+w}{ }\PY{n}{ddof}\PY{o}{=}\PY{l+m+mi}{1}\PY{p}{)}\PY{l+s+si}{:}\PY{l+s+s2}{.2f}\PY{l+s+si}{\PYZcb{}}\PY{l+s+s2}{\PYZdq{}}\PY{p}{)}
\end{Verbatim}
\end{tcolorbox}

    \begin{Verbatim}[commandchars=\\\{\}]
Кластеризация X4 (k-means, k=2)

Кластер 1:
Размер: 114 (57.0\%)
Среднее: 39.30
Медиана: 39.41
Дисперсия: 41.02
Минимум: 22.80
Максимум: 57.44

Кластер 2:
Размер: 86 (43.0\%)
Среднее: 113.71
Медиана: 114.21
Дисперсия: 27.08
Минимум: 100.77
Максимум: 126.84

Общая выборка X4:
Среднее: 71.30
Дисперсия: 1398.94
    \end{Verbatim}

    \begin{tcolorbox}[breakable, size=fbox, boxrule=1pt, pad at break*=1mm,colback=cellbackground, colframe=cellborder]
\prompt{In}{incolor}{24}{\boxspacing}
\begin{Verbatim}[commandchars=\\\{\}]
\PY{c+c1}{\PYZsh{} Визуализация кластеров}

\PY{c+c1}{\PYZsh{} Гистограмма с разделением по кластерам}
\PY{n}{plt}\PY{o}{.}\PY{n}{figure}\PY{p}{(}\PY{n}{figsize}\PY{o}{=}\PY{p}{(}\PY{l+m+mi}{12}\PY{p}{,} \PY{l+m+mi}{5}\PY{p}{)}\PY{p}{)}

\PY{n}{plt}\PY{o}{.}\PY{n}{subplot}\PY{p}{(}\PY{l+m+mi}{1}\PY{p}{,} \PY{l+m+mi}{2}\PY{p}{,} \PY{l+m+mi}{1}\PY{p}{)}
\PY{n}{plt}\PY{o}{.}\PY{n}{hist}\PY{p}{(}\PY{n}{cluster1}\PY{p}{,} \PY{n}{bins}\PY{o}{=}\PY{l+m+mi}{15}\PY{p}{,} \PY{n}{alpha}\PY{o}{=}\PY{l+m+mf}{0.7}\PY{p}{,} \PY{n}{label}\PY{o}{=}\PY{l+s+s1}{\PYZsq{}}\PY{l+s+s1}{Кластер 1}\PY{l+s+s1}{\PYZsq{}}\PY{p}{,} \PY{n}{color}\PY{o}{=}\PY{l+s+s1}{\PYZsq{}}\PY{l+s+s1}{red}\PY{l+s+s1}{\PYZsq{}}\PY{p}{,} \PY{n}{edgecolor}\PY{o}{=}\PY{l+s+s1}{\PYZsq{}}\PY{l+s+s1}{black}\PY{l+s+s1}{\PYZsq{}}\PY{p}{)}
\PY{n}{plt}\PY{o}{.}\PY{n}{hist}\PY{p}{(}\PY{n}{cluster2}\PY{p}{,} \PY{n}{bins}\PY{o}{=}\PY{l+m+mi}{15}\PY{p}{,} \PY{n}{alpha}\PY{o}{=}\PY{l+m+mf}{0.7}\PY{p}{,} \PY{n}{label}\PY{o}{=}\PY{l+s+s1}{\PYZsq{}}\PY{l+s+s1}{Кластер 2}\PY{l+s+s1}{\PYZsq{}}\PY{p}{,} \PY{n}{color}\PY{o}{=}\PY{l+s+s1}{\PYZsq{}}\PY{l+s+s1}{blue}\PY{l+s+s1}{\PYZsq{}}\PY{p}{,} \PY{n}{edgecolor}\PY{o}{=}\PY{l+s+s1}{\PYZsq{}}\PY{l+s+s1}{black}\PY{l+s+s1}{\PYZsq{}}\PY{p}{)}
\PY{n}{plt}\PY{o}{.}\PY{n}{title}\PY{p}{(}\PY{l+s+s1}{\PYZsq{}}\PY{l+s+s1}{Гистограмма X4 с разделением на кластеры}\PY{l+s+s1}{\PYZsq{}}\PY{p}{)}
\PY{n}{plt}\PY{o}{.}\PY{n}{xlabel}\PY{p}{(}\PY{l+s+s1}{\PYZsq{}}\PY{l+s+s1}{Значение}\PY{l+s+s1}{\PYZsq{}}\PY{p}{)}
\PY{n}{plt}\PY{o}{.}\PY{n}{ylabel}\PY{p}{(}\PY{l+s+s1}{\PYZsq{}}\PY{l+s+s1}{Частота}\PY{l+s+s1}{\PYZsq{}}\PY{p}{)}
\PY{n}{plt}\PY{o}{.}\PY{n}{legend}\PY{p}{(}\PY{p}{)}
\PY{n}{plt}\PY{o}{.}\PY{n}{grid}\PY{p}{(}\PY{k+kc}{True}\PY{p}{,} \PY{n}{alpha}\PY{o}{=}\PY{l+m+mf}{0.3}\PY{p}{)}

\PY{c+c1}{\PYZsh{} Boxplot для сравнения}
\PY{n}{plt}\PY{o}{.}\PY{n}{subplot}\PY{p}{(}\PY{l+m+mi}{1}\PY{p}{,} \PY{l+m+mi}{2}\PY{p}{,} \PY{l+m+mi}{2}\PY{p}{)}
\PY{n}{data\PYZus{}to\PYZus{}plot} \PY{o}{=} \PY{p}{[}\PY{n}{cluster1}\PY{p}{,} \PY{n}{cluster2}\PY{p}{,} \PY{n}{X4}\PY{p}{]}
\PY{n}{plt}\PY{o}{.}\PY{n}{boxplot}\PY{p}{(}\PY{n}{data\PYZus{}to\PYZus{}plot}\PY{p}{,} \PY{n}{tick\PYZus{}labels}\PY{o}{=}\PY{p}{[}\PY{l+s+s1}{\PYZsq{}}\PY{l+s+s1}{Кластер 1}\PY{l+s+s1}{\PYZsq{}}\PY{p}{,} \PY{l+s+s1}{\PYZsq{}}\PY{l+s+s1}{Кластер 2}\PY{l+s+s1}{\PYZsq{}}\PY{p}{,} \PY{l+s+s1}{\PYZsq{}}\PY{l+s+s1}{Общая выборка}\PY{l+s+s1}{\PYZsq{}}\PY{p}{]}\PY{p}{)}
\PY{n}{plt}\PY{o}{.}\PY{n}{title}\PY{p}{(}\PY{l+s+s1}{\PYZsq{}}\PY{l+s+s1}{Boxplot сравнения кластеров}\PY{l+s+s1}{\PYZsq{}}\PY{p}{)}
\PY{n}{plt}\PY{o}{.}\PY{n}{ylabel}\PY{p}{(}\PY{l+s+s1}{\PYZsq{}}\PY{l+s+s1}{Значение}\PY{l+s+s1}{\PYZsq{}}\PY{p}{)}
\PY{n}{plt}\PY{o}{.}\PY{n}{grid}\PY{p}{(}\PY{k+kc}{True}\PY{p}{,} \PY{n}{alpha}\PY{o}{=}\PY{l+m+mf}{0.3}\PY{p}{)}

\PY{n}{plt}\PY{o}{.}\PY{n}{tight\PYZus{}layout}\PY{p}{(}\PY{p}{)}
\PY{n}{plt}\PY{o}{.}\PY{n}{show}\PY{p}{(}\PY{p}{)}
\end{Verbatim}
\end{tcolorbox}

    \begin{center}
    \adjustimage{max size={0.9\linewidth}{0.9\paperheight}}{RGR1_A-2_files/RGR1_A-2_57_0.png}
    \end{center}
    { \hspace*{\fill} \\}
    
    Для boxplot

\begin{itemize}
\tightlist
\item
  Нижняя граница ящика (Q1) - 25-й процентиль - 25\% данных меньше этого
  значения
\item
  Линия внутри ящика (Q2) - медиана (50-й процентиль) - центр
  распределения
\item
  Верхняя граница ящика (Q3) - 75-й процентиль - 75\% данных меньше
  этого значения
\item
  Нижний ус - \(Q1 - 1.5 \cdot (Q3 - Q1)\) - нижняя граница для
  невыбросов
\item
  Верхний ус - \(Q3 + 1.5 \cdot (Q3 - Q1)\) - верхняя граница для
  невыбросов
\end{itemize}

    \textbf{Анализ кластеров}

Кластер 1 * Значения сосредоточены в диапазоне 22.80 - 57.44 * Среднее
(39.30) и медиана (39.41) практически совпадают, что указывает на
симметричное распределение внутри кластера * Малая дисперсия (41.02)
свидетельствует о высокой компактности данных * Нет выбросов, так как
значения сосредоточены в узком интервале

Кластер 2 * Значения сосредоточены в диапазоне 100.77 - 126.84 * Среднее
(113.71) и медиана (114.21) близки, распределение внутри кластера также
симметрично * Малая дисперсия (27.08) указывает на компактность данных *
Нет выбросов

Общая выборка: * Объединяет два компактных кластера, расположенных
далеко друг от друга * Дисперсия (1398.94) огромна по сравнению с
дисперсиями внутри кластеров * Среднее (71.30) находится между средними
кластеров и не соответствует центру ни одного из них * Медиана (46.76)
смещена в сторону первого кластера, так как он больше по объему

    \paragraph{Почему общее среднее плохо описывает смесь двух
режимов}\label{ux43fux43eux447ux435ux43cux443-ux43eux431ux449ux435ux435-ux441ux440ux435ux434ux43dux435ux435-ux43fux43bux43eux445ux43e-ux43eux43fux438ux441ux44bux432ux430ux435ux442-ux441ux43cux435ux441ux44c-ux434ux432ux443ux445-ux440ux435ux436ux438ux43cux43eux432}

\begin{enumerate}
\def\labelenumi{\arabic{enumi}.}
\item
  Общее среднее не существует в реальной структуре выборки. Среднее по
  X4 равно 71.30, но в окрестности этого значения почти нет наблюдений:
  между двумя группами есть широкий разрыв.
\item
  Потеря информации о структуре. Общее среднее сглаживает наличие двух
  различных групп. Глядя только на среднее, невозможно понять, что
  данные представляют собой две разные совокупности.
\item
  Ошибочные выводы. Если бы мы использовали только общее среднее, можно
  было бы сказать, что ``типичное'' значение X4 находится около 71.30,
  хотя реально типичных значений два: около 39.30 и около 113.71.
\item
  Маскировка различий. Общее среднее скрывает тот факт, что две группы
  имеют совершенно разные характеристики и расположены далеко друг от
  друга.
\item
  Разница между средними кластеров составляет примерно 74.41 единиц, что
  намного больше внутрикластерного разброса.
\end{enumerate}


    % Add a bibliography block to the postdoc
    
    
    
\end{document}
